% ROZDZIAL 7 -- Lambdy, obiekty funkcyjne i wywolywalne obiekty (C++20)
% Fragment do wlaczenia po rozdziale 6.
% Wymagana preambula: listings, array, amsmath (opcjonalny), hyperref.
% Uzycie: % ROZDZIAL 7 -- Lambdy, obiekty funkcyjne i wywolywalne obiekty (C++20)
% Fragment do wlaczenia po rozdziale 6.
% Wymagana preambula: listings, array, amsmath (opcjonalny), hyperref.
% Uzycie: % ROZDZIAL 7 -- Lambdy, obiekty funkcyjne i wywolywalne obiekty (C++20)
% Fragment do wlaczenia po rozdziale 6.
% Wymagana preambula: listings, array, amsmath (opcjonalny), hyperref.
% Uzycie: % ROZDZIAL 7 -- Lambdy, obiekty funkcyjne i wywolywalne obiekty (C++20)
% Fragment do wlaczenia po rozdziale 6.
% Wymagana preambula: listings, array, amsmath (opcjonalny), hyperref.
% Uzycie: \input{rozdzial_7_lambdy_obiekty_funkcyjne_cpp.tex}

\chapter[Lambdy i obiekty funkcyjne]{Wyrażenia lambda, obiekty funkcyjne i abstrakcje wywoływalne}
\label{chap:lambdas-cpp}

\section{Po co istnieją obiekty funkcyjne?}

W C++ funkcję można reprezentować nie tylko przez jej nazwę lub wskaźnik.
Można również utworzyć \emph{obiekt wywoływalny} (\textit{callable object}):
obiekt, który daje się zastosować w wyrażeniu przypominającym wywołanie
funkcji. Pozwala to przekazywać zachowanie jako parametr algorytmu,
zapamiętywać stan pomiędzy wywołaniami oraz tworzyć elastyczne API
(np. sortowanie z własnym porządkiem, filtrowanie, obsługa zdarzeń).

\begin{lstlisting}
#include <algorithm>
#include <vector>

std::vector<int> numbers{5, 2, 8, 1};
std::sort(numbers.begin(), numbers.end(),
          [](int a, int b) { return a > b; });
// numbers == {8, 5, 2, 1}
\end{lstlisting}

Trzeci argument \texttt{std::sort} to reguła porównywania. Zamiast pisać
osobną funkcję o globalnej nazwie, definiujemy ją dokładnie w miejscu
użycia. Wyrażenie \texttt{[](int a, int b) \{ return a > b; \}}
jest \emph{lambdą}: tworzy obiekt unikatowego typu domknięcia
(\textit{closure type}).

\paragraph{Zakres i wymagania.}
Przyjmujemy \textbf{C++20}. Lambdy wprowadzono w C++11, przechwytywanie
inicjalizujące i generyczne lambdy w C++14, przechwytywanie kopii
\texttt{*this} w C++17, a jawne listy parametrów szablonu lambdy
w C++20. Wykorzystamy semantykę przenoszenia z rozdziału~\ref{chap:move-semantics-cpp},
szablony wariadyczne z rozdziału~\ref{chap:variadic-templates-cpp}
i koncepcje z rozdziału~\ref{chap:concepts-cpp}.

\section{Trzy podstawowe postacie kodu wywoływalnego}
\label{sec:callables-forms}

\subsection{Funkcja i wskaźnik do funkcji}

Zwykła funkcja nie przechowuje kontekstu wywołania. Jej adres można
przypisać wskaźnikowi o odpowiednim typie:

\begin{lstlisting}
int twice(int x) { return 2 * x; }

int (*fp)(int) = &twice;
int a = fp(7);       // 14
int b = (*fp)(7);    // also 14
\end{lstlisting}

W deklaracji \texttt{int (*fp)(int)} nawiasy są istotne. Sam wskaźnik
przechowuje adres funkcji; nie ma miejsca na dodatkowy stan, taki jak
współczynnik mnożenia ustalany w czasie działania programu.

\subsection{Funktor: klasa z \texttt{operator()}}

\begin{lstlisting}
struct Multiply {
    int factor;

    int operator()(int x) const {
        return factor * x;
    }
};

Multiply times_three{3};
int result = times_three(7); // 21
\end{lstlisting}

\texttt{Multiply} jest klasą, a \texttt{times\_three} jej egzemplarzem.
Dzięki polu \texttt{factor} każdy egzemplarz może reprezentować inne
zachowanie. Operator wywołania może być przeciążany, szablonowy,
opatrzony kwalifikatorami \texttt{const} lub \texttt{\&\&} i mieć
złożoną implementację.

\subsection{Lambda jako wygodny zapis funktora}

\begin{lstlisting}
int factor = 3;
auto times_three = [factor](int x) {
    return factor * x;
};
int result = times_three(7); // 21
\end{lstlisting}

Kompilator tworzy odrębny, nienazwany typ klasy odpowiadający lambdzie.
Uproszczony model jej działania jest następujący:

\begin{lstlisting}
struct ClosureModel {
    int factor;

    int operator()(int x) const {
        return factor * x;
    }
};
\end{lstlisting}

Jest to \emph{model pojęciowy}, nie obiecana przez standard reprezentacja
pamięci. Każde wystąpienie wyrażenia lambda ma własny typ, nawet jeżeli
zapis kodu dwóch lambd wygląda identycznie.

\section{Budowa wyrażenia lambda}
\label{sec:lambda-syntax}

Najważniejsze elementy zapisu przedstawia schemat:

\begin{lstlisting}
[captures](parameters) mutable noexcept -> ResultType {
    // body
}
\end{lstlisting}

W C++20 można również jawnie określić parametry szablonu, koncepcje
oraz klauzulę \texttt{requires}. Typ wyniku zwykle jest wywnioskowany
z instrukcji \texttt{return}; podanie \texttt{-> ResultType} bywa
przydatne, gdy chcemy wymusić konkretny typ lub zwracać referencję.

\begin{lstlisting}
auto square = [](int x) { return x * x; };
auto half = [](double x) -> double { return x / 2; };
auto truth = [] { return true; }; // () may be omitted

int a = square(4);  // 16
double b = half(5); // 2.5
\end{lstlisting}

\paragraph{Najważniejsze zastrzeżenie.}
Lambda nie jest wskaźnikiem do funkcji. Zwykle tworzy obiekt z
\texttt{operator()}; jedynie pewne lambdy bez przechwyceń mają
konwersję na wskaźnik do funkcji o pasującej sygnaturze.

\section{Przechwytywanie zmiennych (\textit{captures})}

\begin{center}\small
\begin{tabularx}{\textwidth}{@{}p{0.30\textwidth}X@{}}
\toprule
Zapis lambdy & Znaczenie \\
\midrule
\code{[x]} & Przechwytuje wartość \code{x} przez kopię w obiekcie domknięcia. \\
\texttt{[\&x]} & Przechwytuje referencję do \code{x}; wywołanie wymaga, by \code{x} nadal żyło. \\
\code{[=]} & Domyślnie przechwytuje użyte zmienne lokalne przez wartość. \\
\texttt{[\&]} & Domyślnie przechwytuje użyte zmienne lokalne przez referencję. \\
\code{[p = std::move(p)]} & Przenosi wartość do składowej domknięcia przez inicjalizację przechwycenia. \\
\code{mutable} & Pozwala zmieniać kopie przechwycone przez wartość wewnątrz lambdy. \\
\bottomrule
\end{tabularx}
\end{center}
\label{sec:lambda-captures}

\subsection{Przechwycenie przez wartość}

\begin{lstlisting}
int threshold = 10;
auto above = [threshold](int x) {
    return x > threshold;
};
auto above_current = [&threshold](int x) {
    return x > threshold;
};
threshold = 100;

bool first = above(20);         // true: zapamietane 10
bool second = above_current(20); // false: aktualne 100
\end{lstlisting}

\texttt{[threshold]} inicjalizuje składową obiektu domknięcia wartością
\texttt{threshold} w chwili tworzenia lambdy. Późniejsza zmiana zmiennej
zewnętrznej nie modyfikuje tej kopii. Jeśli kopiowany obiekt ma kosztowne
kopiowanie, samo utworzenie lub skopiowanie lambdy może także mieć koszt.

\subsection{Przechwycenie przez referencję}

\begin{lstlisting}
int total = 0;
auto add = [&total](int value) {
    total += value;
};
add(4);
add(8); // total == 12
\end{lstlisting}

\texttt{[\&total]} uzyskuje dostęp do tej samej zmiennej, nie do jej
kopii. Lambda \textbf{nie wydłuża czasu życia} obiektu przechwyconego
przez referencję. Dlatego jest wygodna, gdy callback wykonywany jest
synchronicznie, ale potencjalnie niebezpieczna, jeśli jest zapamiętywany
albo wykonywany później.

\subsection{Przechwycenia domyślne i mieszane}

\begin{lstlisting}
int a = 2;
int b = 3;
int count = 0;

auto first  = [=]  { return a + b; }; // by copy
auto second = [&]  { ++count; };     // by reference
auto third  = [=, &count] {
    ++count;
    return a + b;
};
\end{lstlisting}

\texttt{[=]} i \texttt{[\&]} obejmują te użyte w treści zmienne lokalne,
które podlegają niejawnemu przechwyceniu; nie oznaczają automatycznego
skopiowania \emph{wszystkich} zmiennych z otoczenia. W większym kodzie
zwykle lepiej wypisać przechwycenia jawnie: dzięki temu widać zależności
i łatwiej ocenić bezpieczeństwo czasu życia.

\subsection{Przechwytywanie inicjalizujące (C++14)}

W nawiasach kwadratowych można zdefiniować nową nazwę i sposób jej
inicjalizacji. Dzięki temu lambdy mogą przechowywać wartości pochodne,
a także zasoby, których nie wolno kopiować:

\begin{lstlisting}
#include <memory>
#include <utility>

auto p = std::make_unique<int>(42);
auto read = [q = std::move(p)] {
    return *q;
};
// p == nullptr; ownership is now in the closure
int answer = read(); // 42
\end{lstlisting}

Taka lambda jest \emph{move-only}, ponieważ zawiera
\texttt{std::unique\_ptr}. Nie można jej skopiować. To istotne przy
wyborze miejsca przechowywania callbacków: w C++20
\texttt{std::function} wymaga kopiowalnego obiektu docelowego.

\subsection{Niebezpieczne przechwytywanie: wisząca referencja}

\begin{lstlisting}
// WRONG: do not call the returned lambda!
auto broken() {
    int local = 5;
    return [&local] { return local; };
}

// Correct: capture the value.
auto safe() {
    int local = 5;
    return [local] { return local; };
}
\end{lstlisting}

\texttt{broken()} zwraca obiekt z dostępem do zmiennej, która przestaje
istnieć wraz z wyjściem z funkcji. Późniejsze wywołanie prowadziłoby
do niezdefiniowanego zachowania (\textit{undefined behavior}).
Przechwycenie przez kopię rozwiązuje ten konkretny problem.

\section{Przechwytywanie \texttt{this} i \texttt{*this}}
\label{sec:lambda-this}

W metodzie klasy można korzystać z jej pól wewnątrz lambdy. Trzeba
jednak rozróżnić przechwycenie wskaźnika \texttt{this} od skopiowania
bieżącego obiektu.

\begin{lstlisting}
#include <string>

struct Label {
    std::string text;

    auto by_pointer() const {
        return [this] { return text; };
    }

    auto by_copy() const {
        return [*this] { return text; }; // C++17
    }
};
\end{lstlisting}

Lambda zwrócona przez \texttt{by\_pointer()} odwołuje się do oryginalnego
obiektu. Nie może być bezpiecznie wywoływana po jego zniszczeniu.
Lambda z \texttt{[*this]} przechowuje kopię obiektu (jeśli da się go
skopiować); dla składowych zarządzanych przez RAII jest to zwykle
bardziej przewidywalne, ale może być kosztowne.

\textbf{Uwaga:} zapis \texttt{[=]} użyty w metodzie nie kopiuje automatycznie
całego obiektu. Niejawne przechwytywanie \texttt{this} przez \texttt{[=]}
jest w C++20 przestarzałe (\textit{deprecated}); zapisujmy
\texttt{[this]} jawnie albo \texttt{[*this]}, gdy potrzebujemy kopii.

\section{Domyślnie stały operator wywołania i \texttt{mutable}}

Dla typowej lambdy bez słowa \texttt{mutable} operator wywołania ma
kwalifikator \texttt{const}. Oznacza to, że przechwyconych przez wartość
składowych nie wolno w nim zwyczajnie modyfikować:

\begin{lstlisting}
int counter = 0;

// ERROR: captured counter is not writable in this closure.
// auto bad = [counter] { ++counter; };

auto next = [counter]() mutable {
    return ++counter;
};

int a = next(); // 1
int b = next(); // 2
// outer counter is still 0
\end{lstlisting}

\texttt{mutable} modyfikuje zachowanie \emph{kopii wewnątrz lambdy},
nie powoduje przechwycenia oryginalnej zmiennej przez referencję.
Różne kopie obiektu domknięcia mogą przechowywać różny stan:

\begin{lstlisting}
auto a = [n = 0]() mutable { return ++n; };
auto b = a; // independent copy of state

a(); // 1
a(); // 2
b(); // 1
\end{lstlisting}

Przy przechwyceniu przez referencję do modyfikacji zewnętrznego
obiektu \texttt{mutable} nie jest potrzebne: stałość metody
\texttt{operator()} nie zmienia stałości obiektu osiąganego przez
referencję.

\section{Typ lambdy, kopiowanie i konwersja na wskaźnik}

\begin{lstlisting}
#include <type_traits>

inline constexpr auto f = [](int x) { return x + 1; };

auto a = [](int x) { return x; };
auto b = [](int x) { return x; };
static_assert(!std::is_same_v<decltype(a), decltype(b)>);

int (*pf)(int) = +[](int x) { return x + 1; };
\end{lstlisting}

Dwie osobno zapisane lambdy mają różne typy, nawet jeśli robią to samo.
W C++20 typ domknięcia lambdy bez przechwyceń jest domyślnie
konstruowalny; dostępność różnych operacji kopiowania/przypisywania
zależy od tego, czy lambda przechwytuje zmienne. Dla lambd z
przechwyceniami obowiązują też ograniczenia narzucone przez typy pól
(na przykład niekopiowalność \texttt{unique\_ptr}).

\texttt{+[]} jest znanym idiomem wymuszającym konwersję
nieprzechwytującej żadnej zmiennej lambdy na wskaźnik do funkcji.
Nie działa dla lambd przechwytujących:

\begin{lstlisting}
int offset = 10;
auto with_state = [offset](int x) { return x + offset; };
// int (*bad)(int) = with_state;  // ERROR
\end{lstlisting}

\section{Lambdy generyczne i własne parametry szablonu}
\label{sec:lambda-generic}

\subsection{Parametry \texttt{auto} (C++14)}

\begin{lstlisting}
#include <string>
auto print_twice = [](const auto& value) {
    return value + value;
};

auto x = print_twice(7);                 // 14
auto s = print_twice(std::string("ab")); // "abab"
\end{lstlisting}

Parametr \texttt{auto} sprawia, że \texttt{operator()} obiektu domknięcia
jest szablonem funkcji. Przy wywołaniu kompilator dedukuje odpowiednie
typy tak jak przy zwykłej funkcji szablonowej.

\subsection{Jawne parametry szablonu w lambdzie (C++20)}

\begin{lstlisting}
#include <concepts>

auto equal_types = []<class T>(const T& a, const T& b) {
    return a == b;
};

auto sum_integrals = []<std::integral... Ts>(Ts... values) {
    return (0 + ... + values);
};

bool yes = equal_types(3, 3);
int total = sum_integrals(1, 2, 3, 4); // 10
\end{lstlisting}

W przypadku \texttt{equal\_types} oba argumenty muszą dać ten sam
dedukowany typ \texttt{T}. W drugim przykładzie łączymy \emph{concepts}
z pakietem parametrów i wyrażeniem fold znanym z
rozdziałów~\ref{chap:variadic-templates-cpp} i~\ref{chap:concepts-cpp}.

\subsection{Uniwersalny wrapper z perfect forwarding}

\begin{lstlisting}
#include <functional>
#include <utility>

auto invoke_once = []<class F, class... Args>(
    F&& f, Args&&... args) -> decltype(auto) {
    return std::invoke(std::forward<F>(f),
                       std::forward<Args>(args)...);
};
\end{lstlisting}

\texttt{F\&\&} i \texttt{Args\&\&...} są tu referencjami przekazującymi,
a \texttt{decltype(auto)} zachowuje typ wyniku, w tym jego
referencyjność. W odróżnieniu od \texttt{std::move},
\texttt{std::forward} nie wymusza kategorii rvalue dla każdego argumentu.
Dzięki \texttt{std::invoke} wrapper działa nie tylko dla zwykłych
funkcji i lambd, lecz także dla wskaźników do składowych.

\section{Wymuszanie sygnatury: \texttt{noexcept}, \texttt{requires}, zwracanie referencji}

\begin{lstlisting}
#include <concepts>
#include <utility>

auto positive = [](int x) noexcept {
    return x > 0;
};

auto increment = []<std::integral T>(T x) -> T {
    return x + 1;
};

auto first = []<class T>(T& x) -> T& {
    return x;
};

int value = 3;
first(value) = 17; // modifies value
\end{lstlisting}

Ważne są dwie subtelności. Po pierwsze, \texttt{noexcept} jest obietnicą
programisty; jeśli wyjątek wydostanie się z funkcji oznaczonej
\texttt{noexcept}, wywołane zostanie \texttt{std::terminate}.
Po drugie, domyślna dedukcja wyniku lambdy deklarowanej bez
\texttt{decltype(auto)} lub jawnego \texttt{-> T\&} zwykle usuwa
referencję. Właściwe zwracanie referencji wymaga zatem świadomego
określenia typu wyniku i zapewnienia jego poprawnego czasu życia.

\section{\texttt{std::invoke}: jeden mechanizm wywołania}

\begin{center}\small
\begin{tabularx}{\textwidth}{@{}p{0.34\textwidth}X@{}}
\toprule
Narzędzie & Za co odpowiada \\
\midrule
\code{std::invoke(f, args...)} & Jednolicie wywołuje funkcję, lambdę lub wskaźnik do składowej. \\
\code{std::function<R(Args...)>} & Przechowuje kopiowalny obiekt wywoływalny o ustalonej sygnaturze przez wymazywanie typu. \\
\code{std::ref(x)} & Opakowuje referencję, gdy API normalnie kopiuje argument. \\
\code{std::bind_front(f, args...)} & Wiąże pierwsze argumenty funkcji i zwraca nowy obiekt wywoływalny. \\
\code{noexcept} przy lambdzie & Określa, że operator wywołania nie powinien pozwolić wyjątkowi opuścić funkcji. \\
\bottomrule
\end{tabularx}
\end{center}
\label{sec:std-invoke}

\texttt{std::invoke} z nagłówka \texttt{<functional>} wywołuje funkcję,
funktor lub lambdę, ale obsługuje także wskaźnik do metody i pola.

\begin{lstlisting}
#include <functional>

struct Box {
    int value = 10;
    int add(int x) const { return value + x; }
};

Box b;
int a = std::invoke([](int x) { return x * 2; }, 4);
int c = std::invoke(&Box::add, b, 7);      // 17
int& r = std::invoke(&Box::value, b);
r = 25;                                   // b.value == 25
\end{lstlisting}

To podstawowa cegiełka generycznych algorytmów. W nagłówku
\texttt{<type\_traits>} zdefiniowano m.in. \texttt{std::invoke\_result\_t},
\texttt{std::is\_invocable\_v} i \texttt{std::is\_nothrow\_invocable\_v},
które pozwalają sprawdzić w czasie kompilacji poprawność wywołania.

\begin{lstlisting}
#include <type_traits>

static_assert(std::is_invocable_v<decltype(positive), int>);
static_assert(!std::is_invocable_v<decltype(positive),
                                   std::string>);
static_assert(std::is_nothrow_invocable_v<decltype(positive),
                                          int>);
\end{lstlisting}

\section{\texttt{std::function}: jeden typ dla wielu zachowań}
\label{sec:std-function}

Lambdy o różnych typach nie mogą bezpośrednio zastępować się w jednej
zmiennej zadeklarowanej przez \texttt{auto}. Jeżeli chcemy przechowywać
\emph{dowolny} kopiowalny obiekt wywoływalny o zadanej sygnaturze,
możemy użyć \texttt{std::function<R(Args...)>}.

\begin{lstlisting}
#include <functional>
#include <string>

std::function<int(int)> operation =
    [](int x) { return 2 * x; };

int factor = 5;
operation = [factor](int x) { return factor * x; };
int result = operation(3); // 15

operation = nullptr;
if (!operation) {
    // no callable stored
}
\end{lstlisting}

\texttt{std::function} stosuje \emph{type erasure}: ukrywa rzeczywisty
typ obiektu docelowego za wspólnym interfejsem wywołania. Może
przechowywać funkcje, funktory i lambdy z przechwyceniami. Nie jest
odpowiednikiem \texttt{auto}: deklaracja \texttt{auto} zachowuje
konkretny typ i pozwala zwykle kompilatorowi łatwiej zoptymalizować
wywołanie.

\subsection{Pusty \texttt{std::function} i wyjątek}

\begin{lstlisting}
#include <functional>

std::function<void()> callback;
if (callback) {
    callback();
}
// Calling callback() while empty throws std::bad_function_call.
\end{lstlisting}

Sprawdzenie \texttt{if (callback)} jest zatem standardową częścią
obsługi opcjonalnych callbacków. Nie należy przechowywać
\texttt{std::function} bez potrzeby, jeśli konkretny typ wywoływalny
jest znany w czasie kompilacji.

\subsection{Koszty i ograniczenie kopiowalności}

Zwykle \texttt{std::function} wymaga pośredniego wywołania. Przechowywanie
większego domknięcia może też wymagać dynamicznej alokacji (zależy to
od implementacji i rozmiaru obiektu). To koszt elastyczności API, nie
argument przeciw używaniu \texttt{std::function} w każdym przypadku.

W C++20 przechowywany obiekt docelowy musi być kopiowalny, dlatego
następujący kod nie jest poprawny:

\begin{lstlisting}
#include <functional>
#include <memory>

// ERROR in C++20: the closure is not copy-constructible.
// std::function<int()> f =
//     [p = std::make_unique<int>(9)] { return *p; };
\end{lstlisting}

Od \textbf{C++23} biblioteka udostępnia
\texttt{std::move\_only\_function}, która służy do przechowywania
nienadających się do kopiowania obiektów wywoływalnych. W C++20
można natomiast przechowywać taki obiekt bezpośrednio jako \texttt{auto}
lub przekazać go do szablonu, który nie wymaga kopiowania.

\section{Co wybrać: \texttt{auto}, szablon, wskaźnik czy \texttt{std::function}?}
\label{sec:lambda-choice}

\begin{center}
\begin{tabular}{|p{3.0cm}|p{4.4cm}|p{5.0cm}|}
\hline
\textbf{Forma} & \textbf{Zalety} & \textbf{Ograniczenia} \\
\hline
\texttt{auto f = lambda} & Konkretny typ, bez konieczności wymazywania typu, naturalna obsługa stanu & Zmienna ma ustalony typ domknięcia \\
\hline
\texttt{template<class F>} & Działa z wieloma typami, umożliwia inline i forwarding & Definicja szablonu zwykle dostępna w nagłówku; możliwe powiększanie kodu \\
\hline
Wskaźnik do funkcji & Prosty nieposiadający stanu interfejs & Nie przyjmuje lambd z przechwyceniami \\
\hline
\texttt{std::function} & Jedna sygnatura, zmienne zachowanie w czasie wykonania & Type erasure; koszt pośredniego wywołania; wymóg kopiowalności celu w C++20 \\
\hline
\end{tabular}
\end{center}

\paragraph{Praktyczna reguła.}
Jeżeli callback wykorzystujemy natychmiast, wybierajmy zwykle szablon
lub \texttt{auto}. Jeżeli musimy przechowywać i później podmieniać różne
callbacki o jednej sygnaturze, \texttt{std::function} jest często
czytelniejsza niż ręczne modelowanie polimorfizmu.

\section{Lambdy ze standardowymi algorytmami}
\label{sec:lambda-algorithms}

\subsection{\texttt{find\_if}, \texttt{count\_if}, \texttt{transform}}

\begin{lstlisting}
#include <algorithm>
#include <vector>

std::vector<int> xs{1, 2, 3, 4, 5, 6};

int limit = 3;
auto it = std::find_if(xs.begin(), xs.end(),
                       [limit](int x) { return x > limit; });
// *it == 4

auto even_count = std::count_if(xs.begin(), xs.end(),
                                [](int x) { return x % 2 == 0; });
// even_count == 3

std::vector<int> squares(xs.size());
std::transform(xs.begin(), xs.end(), squares.begin(),
               [](int x) { return x * x; });
\end{lstlisting}

Dla \texttt{std::find\_if} trzeba przed dereferencją sprawdzić, czy
wynik jest różny od \texttt{xs.end()}. Dla \texttt{std::transform} zbiór
docelowy musi już zawierać wystarczającą liczbę elementów: samo
\texttt{reserve()} zwiększa \texttt{capacity()}, lecz nie \texttt{size()}.
Alternatywą jest pusty wektor i iterator \texttt{std::back\_inserter}.

\subsection{\texttt{std::ranges} i projekcje (C++20)}

\begin{lstlisting}
#include <algorithm>
#include <ranges>
#include <string>
#include <vector>

struct Person {
    std::string name;
    int age;
};

std::vector<Person> people{{"Ada", 31}, {"Bob", 18}, {"Eve", 25}};

std::ranges::sort(people, {}, &Person::age);
// sorted by age ascending

auto adults = people | std::views::filter(
    [](const Person& p) { return p.age >= 18; });
\end{lstlisting}

Trzeci argument \texttt{ranges::sort} to \emph{projekcja}:
\texttt{\&Person::age} pobiera klucz sortowania, a domyślny komparator
porównuje klucze. \texttt{views::filter} nie buduje automatycznie nowego
wektora; jest leniwym widokiem, który odwołuje się do danych źródłowych.
Przy widokach szczególnie uważnie trzeba kontrolować czas życia
kontenera i przechwyceń predykatu.

\section{\texttt{std::ref}, \texttt{std::bind\_front} i częściowe wiązanie argumentów}
\label{sec:bind-front}

\begin{lstlisting}
#include <functional>

int add(int x, int y) { return x + y; }

auto add_ten = std::bind_front(add, 10); // C++20
int x = add_ten(5);                      // 15
\end{lstlisting}

\texttt{std::bind\_front} przechowuje funkcję i ustalone argumenty,
a pozostałe przyjmuje podczas wywołania. W wielu zastosowaniach
równoważna lambda jest czytelniejsza:

\begin{lstlisting}
auto add_ten_lambda = [](int x) { return add(10, x); };
\end{lstlisting}

Domyślnie argumenty przechowywane przez \texttt{bind\_front} są
przechowywane przez wartość, z uwzględnieniem przenoszenia. Jeśli
potrzebujemy odwołania do oryginalnego obiektu, stosujemy
\texttt{std::ref} (a dla dostępu tylko do odczytu
\texttt{std::cref}), pamiętając o czasie jego życia:

\begin{lstlisting}
#include <functional>

struct Accumulator {
    int total = 0;
    void add(int n) { total += n; }
};

Accumulator acc;
auto update = std::bind_front(&Accumulator::add, std::ref(acc));
update(7); // acc.total == 7
\end{lstlisting}

Starsze \texttt{std::bind} jest dostępne, ale ma bardziej skomplikowane
reguły dotyczące argumentów i podstawiania. W nowym kodzie zwykle
preferujemy lambdy lub \texttt{std::bind\_front}.

\section{Praktyczny projekt: rejestr callbacków}
\label{sec:callbacks}

Rozważmy klasę powiadamiającą o zmianie wartości. Do przechowywania
procedury obsługi zdarzenia użyjemy \texttt{std::function}. Dzięki
temu można ją zastąpić inną nawet po utworzeniu obiektu.

\begin{lstlisting}
#include <functional>
#include <utility>

class Counter {
public:
    using Callback = std::function<void(int)>;

    void on_change(Callback callback) {
        callback_ = std::move(callback);
    }

    void increment() {
        ++value_;
        if (callback_) callback_(value_);
    }

    int value() const noexcept { return value_; }

private:
    int value_ = 0;
    Callback callback_;
};

Counter counter;
counter.on_change([](int value) {
    // report the new value
});
counter.increment();
\end{lstlisting}

Kod w pojedynczym wątku jest prosty i wystarczający dla wielu
interfejsów zdarzeniowych. W programie wielowątkowym równoczesna
zmiana \texttt{callback\_} oraz wykonywanie \texttt{increment()}
wymaga dodatkowej synchronizacji. \texttt{std::function} nie zapewnia
automatycznie ochrony wewnętrznego stanu przed wyścigami danych.

\subsection{Callback a czas życia właściciela}

Typowy przebieg to: rejestrujemy callback, właściciel danych zostaje
zniszczony, a ktoś później wywołuje callback. Jeśli lambda zachowała
samą referencję lub surowy wskaźnik, dostęp byłby wtedy błędny.
Słaby wskaźnik w poniższym przykładzie pozwala sprawdzić czas życia.

\begin{lstlisting}
#include <memory>

struct Worker : std::enable_shared_from_this<Worker> {
    int total = 0;

    auto callback() {
        std::weak_ptr<Worker> weak = shared_from_this();
        return [weak](int n) {
            if (auto self = weak.lock()) {
                self->total += n;
            }
        };
    }
};
\end{lstlisting}

Przykład zakłada, że obiekt \texttt{Worker} jest już zarządzany przez
\texttt{std::shared\_ptr} w chwili wywołania \texttt{callback()}.
\texttt{weak\_ptr} pozwala sprawdzić, czy obiekt wciąż istnieje,
bez przedłużania jego życia przez przechwycenie silnego wskaźnika.
To ważne również dla unikania cykli własności, np. gdy sam obiekt
przechowuje callback odnoszący się do niego.

\section{Rekurencyjne lambdy bez \texttt{std::function}}

Lambda zwykle nie zna bezpośrednio własnej nazwy podczas definiowania.
Można jednak przekazywać obiekt wywoływalny jako argument:

\begin{lstlisting}
auto factorial = [](auto&& self, unsigned n) -> unsigned {
    return n <= 1 ? 1 : n * self(self, n - 1);
};

unsigned x = factorial(factorial, 5); // 120
\end{lstlisting}

To odmiana wzorca \emph{self-parameter}. Dla realnych zastosowań
rekurencyjnych warto ponadto zwrócić uwagę na przepełnienie liczb
całkowitych i maksymalną głębokość stosu. W C++20 metoda ta pozwala
uniknąć pośredniej alokacji i wymazywania typu, których mogłaby wymagać
rekurencja oparta na \texttt{std::function}.

\section{Lambdy w metaprogramowaniu: \texttt{std::visit} i \texttt{std::apply}}
\label{sec:lambda-visit-apply}

Dwa przydatne zastosowania łączą lambdy z mechanizmami z
rozdziałów~\ref{chap:cpp14-type-erasure} i~\ref{chap:variadic-templates-cpp}.
\texttt{std::visit} wybiera wywołanie na podstawie aktywnego typu
\texttt{std::variant}, natomiast \texttt{std::apply} rozwija argumenty
z krotki lub podobnej struktury.

\subsection{Zestaw przeciążeń oparty na dziedziczeniu}

\begin{lstlisting}
#include <string>
#include <variant>

template <class... Fs>
struct overloaded : Fs... {
    using Fs::operator()...;
};

std::variant<int, std::string> value = std::string{"Hello"};

auto stringify = overloaded{
    [](int x) { return std::to_string(x); },
    [](const std::string& x) { return x; }
};

std::string result = std::visit(stringify, value); // "Hello"
\end{lstlisting}

Każda lambda dostarcza jeden \texttt{operator()}. Dziedziczenie
wariadyczne i deklaracja \texttt{using Fs::operator()...} tworzą zbiór
przeciążeń, spośród których kompilator wybiera odpowiednią funkcję.
W tym przykładzie wszystkie warianty wywołania zwracają ten sam typ,
\texttt{std::string}, czego wymaga tu pojedyncze wywołanie
\texttt{std::visit}. Konstrukcja jest przydatna przy przetwarzaniu
stanów programu i heterogenicznych danych.

\subsection{Rozwijanie krotki przez \texttt{std::apply}}

\begin{lstlisting}
#include <tuple>

std::tuple<int, int, int> data{2, 3, 5};
int sum = std::apply(
    [](auto... values) { return (0 + ... + values); },
    data);
// sum == 10
\end{lstlisting}

Tutaj \texttt{std::apply} bierze elementy krotki i wywołuje lambdę
jak zwykłą funkcję trójargumentową. Sama lambda wykorzystuje
\emph{fold expression}, a więc nie trzeba pisać rekurencyjnego kodu.

\section{Pakiety przechwyceń i obliczenia odroczone (C++20)}

C++20 pozwala stosować \emph{pack expansion} również w przechwytywaniu
inicjalizującym. Możemy w ten sposób zbudować jeden obiekt domknięcia,
który przechowuje cały pakiet wartości do późniejszego przetworzenia:

\begin{lstlisting}
#include <concepts>

template <std::integral... Ts>
auto delayed_sum(Ts... xs) {
    return [...stored = xs] {
        return (0 + ... + stored);
    };
}

auto job = delayed_sum(2, 3, 5);
int result = job(); // 10
\end{lstlisting}

Lambdę można zwrócić z funkcji, ponieważ wartości są przechowywane
wewnątrz domknięcia. Przechwycenie przez referencję miałoby inne
wymagania dotyczące czasu życia. W rzeczywistych fabrykach callbacków
pakiet może zawierać obiekty niekopiowalne; wtedy należy zadbać
o przekazanie własności (np. przez \texttt{std::move} i
\texttt{std::forward}) oraz o wymagania interfejsu, który callback
przechowuje.

\section{Lambdy w obliczeniach stałych}

Jeśli ciało lambdy dopuszcza obliczanie stałe, jej wywołanie może
być częścią wyrażenia stałego. Od C++17 odpowiedni operator wywołania
jest domyślnie \texttt{constexpr}, gdy spełnia wymagania takiej funkcji:

\begin{lstlisting}
constexpr auto cube = [](int x) { return x * x * x; };
static_assert(cube(3) == 27);

// C++20: each call must be evaluated at compile time.
auto fourth_power = [](int x) consteval {
    return x * x * x * x;
};
static_assert(fourth_power(2) == 16);
\end{lstlisting}

\texttt{consteval} jest mocniejszą deklaracją niż \texttt{constexpr}:
wymaga, aby właściwe wywołania funkcji natychmiastowej spełniały
reguły obliczania stałego. Nie należy traktować tego jako gwarancji,
że każda zwykła lambda zostanie wykonana podczas kompilacji.

\section{Koncepcje dla obiektów wywoływalnych}
\label{sec:callable-concepts}

W C++20 można ograniczać typy wywoływalne m.in. następującymi
koncepcjami:
\begin{itemize}
\item \texttt{std::invocable} -- poprawność wywołania z danymi argumentami;
\item \texttt{std::regular\_invocable} -- dodatkowe wymagania semantyczne;
\item \texttt{std::predicate} -- poprawnie wywoływalny predykat z
wynikiem dającym się traktować jako wartość logiczną.
\end{itemize}
Koncepcje nie dowodzą wszystkich własności semantycznych. Przykładowo
kompilator nie zweryfikuje matematycznie stabilności wyników predykatu.

\begin{lstlisting}
#include <concepts>
#include <functional>
#include <utility>

template <class F, class T>
    requires std::predicate<F&, const T&>
bool test(const T& value, F&& predicate) {
    return std::invoke(std::forward<F>(predicate), value);
}

bool ok = test(8, [](int x) { return x % 2 == 0; });
\end{lstlisting}

\textbf{Istotne rozróżnienie:} koncepcja weryfikuje przede wszystkim
własności typów i poprawność wyrażeń. Nie dowodzi matematycznie
własności predykatu, np. że comparator tworzy ścisły słaby porządek.
Programista nadal odpowiada za kontrakt semantyczny.

\section{Typowe błędy i trudniejsze przypadki}
\label{sec:lambda-pitfalls}

\subsection{Referencja przechwycona do operacji asynchronicznej}

\begin{lstlisting}
#include <future>

std::future<int> start() {
    int x = 7;
    // WRONG: the task could outlive x.
    // return std::async(std::launch::async, [&x] { return x; });
    return std::async(std::launch::async, [x] { return x; });
}
\end{lstlisting}

W zadaniach asynchronicznych preferuj przechwycenia przez wartość
lub świadomie zarządzaną własność, dopóki nie mamy ścisłej gwarancji,
że obiekt referowany przetrwa wykonanie zadania. Sama biblioteka
\texttt{std::async} nie wydłuża życia lokalnych zmiennych funkcji
uruchamiającej.

\subsection{Predykat sortowania nie może naruszać porządku}

\begin{lstlisting}
// WRONG: <= is not a strict ordering relation.
// std::sort(xs.begin(), xs.end(),
//           [](int a, int b) { return a <= b; });

std::sort(xs.begin(), xs.end(),
          [](int a, int b) { return a < b; });
\end{lstlisting}

Komparator do \texttt{sort} musi definiować \emph{ścisły słaby
porządek}. Przykładowo \texttt{comp(x,x)} ma dawać \texttt{false}.
Naruszenie wymagań algorytmu może prowadzić do niezdefiniowanego
zachowania, mimo że typy argumentów się zgadzają.

\subsection{Pomylenie kopii z oryginalnym stanem}

\begin{lstlisting}
int sum = 0;
std::for_each(xs.begin(), xs.end(), [sum](int x) mutable {
    sum += x; // modifies only the closure's copy
});
// outer sum remains 0

std::for_each(xs.begin(), xs.end(), [&sum](int x) {
    sum += x;
});
\end{lstlisting}

Nie należy zakładać, że \texttt{mutable} zmienia zewnętrzną
zmienną. Należy też pamiętać, że algorytmy
mogą kopiować obiekt funkcyjny; nie opierajmy semantyki algorytmu
na nieudokumentowanych szczegółach kopiowania predykatu.

\subsection{Przechwycenie wskaźnika nie przedłuża życia obiektu}

\begin{lstlisting}
int* p = /* pointer to some object */ nullptr;
auto f = [p] {
    if (p) return *p;
    return 0;
};
\end{lstlisting}

Wartość wskaźnika może zostać skopiowana do lambdy, ale wskazywany
obiekt nadal musi istnieć w chwili dereferencji. Przechwycenie
\texttt{shared\_ptr} przenosi własność współdzieloną, a
\texttt{weak\_ptr} pozwala ją warunkowo odzyskać. To nie jest to
samo co skopiowanie wskaźnika surowego.

\section{Zestawienie: kiedy stosować którą konstrukcję?}

\begin{center}
\begin{tabular}{|p{5.7cm}|p{6.6cm}|}
\hline
\textbf{Potrzeba} & \textbf{Zalecany wybór} \\
\hline
Krótki predykat do algorytmu & Lambda z jawnymi przechwyceniami \\
\hline
Wywołanie bez stanu w API C & Zwykła funkcja albo zgodna lambda bez przechwyceń jako wskaźnik do funkcji \\
\hline
Wydajny callback o typie znanym w czasie kompilacji & Parametr szablonowy \texttt{F}, ewentualnie \texttt{std::invocable} \\
\hline
Callback z możliwością późniejszej podmiany & \texttt{std::function} (C++20, kopiowalny cel) \\
\hline
Przechowanie \texttt{unique\_ptr} wewnątrz lambdy & Init-capture i przeniesienie; nie przechowuj takiej lambdy w \texttt{std::function} w C++20 \\
\hline
Callback żyjący dłużej niż zmienne lokalne & Kopiuj potrzebny stan lub zarządzaj życiem przez RAII i smart pointery \\
\hline
Uniwersalne wywołanie dowolnego callable & \texttt{std::invoke} i \texttt{std::forward} \\
\hline
Częściowe podstawienie kilku argumentów & Lambda lub \texttt{std::bind\_front} \\
\hline
Bezpieczny predykat dla szablonu & \texttt{std::predicate} plus kontrola jego semantyki \\
\hline
\end{tabular}
\end{center}

\paragraph{Pytania kontrolne.}
\begin{enumerate}
\item Dlaczego dwie lambdy o identycznym kodzie mają różne typy?
\item Jaka jest różnica między \texttt{[x]}, \texttt{[\&x]} i \texttt{[x]() mutable}?
\item Kiedy bezpiecznie zwrócić lambdę przechwytującą zmienną lokalną?
\item Dlaczego lambda posiadająca \texttt{unique\_ptr} nie może być celem \texttt{std::function} w C++20?
\item Jak \texttt{std::invoke} wywołuje wskaźnik do metody klasy?
\item Czym różnią się \texttt{[this]} i \texttt{[*this]}?
\item Co może pójść źle, gdy \texttt{std::views::filter} korzysta z lokalnej zmiennej przez referencję?
\item Kiedy interfejs przyjmujący szablon \texttt{F} jest lepszy od \texttt{std::function}?
\item Dlaczego komparator \texttt{a <= b} nie jest poprawny dla \texttt{std::sort}?
\item W jakiej sytuacji przechwycenie \texttt{weak\_ptr} jest lepsze od \texttt{shared\_ptr}?
\end{enumerate}

\paragraph{Dalsza lektura.}
W dokumentacji \emph{cppreference} warto przeczytać hasła
\emph{Lambda expressions}, \emph{Function objects} i \emph{Ranges}.
Dokładne opisy narzędzi znajdują się również pod hasłami:
\begin{itemize}
\item \texttt{std::invoke} oraz \texttt{std::function};
\item \texttt{std::bind\_front} oraz \texttt{std::predicate}.
\end{itemize}


\chapter[Lambdy i obiekty funkcyjne]{Wyrażenia lambda, obiekty funkcyjne i abstrakcje wywoływalne}
\label{chap:lambdas-cpp}

\section{Po co istnieją obiekty funkcyjne?}

W C++ funkcję można reprezentować nie tylko przez jej nazwę lub wskaźnik.
Można również utworzyć \emph{obiekt wywoływalny} (\textit{callable object}):
obiekt, który daje się zastosować w wyrażeniu przypominającym wywołanie
funkcji. Pozwala to przekazywać zachowanie jako parametr algorytmu,
zapamiętywać stan pomiędzy wywołaniami oraz tworzyć elastyczne API
(np. sortowanie z własnym porządkiem, filtrowanie, obsługa zdarzeń).

\begin{lstlisting}
#include <algorithm>
#include <vector>

std::vector<int> numbers{5, 2, 8, 1};
std::sort(numbers.begin(), numbers.end(),
          [](int a, int b) { return a > b; });
// numbers == {8, 5, 2, 1}
\end{lstlisting}

Trzeci argument \texttt{std::sort} to reguła porównywania. Zamiast pisać
osobną funkcję o globalnej nazwie, definiujemy ją dokładnie w miejscu
użycia. Wyrażenie \texttt{[](int a, int b) \{ return a > b; \}}
jest \emph{lambdą}: tworzy obiekt unikatowego typu domknięcia
(\textit{closure type}).

\paragraph{Zakres i wymagania.}
Przyjmujemy \textbf{C++20}. Lambdy wprowadzono w C++11, przechwytywanie
inicjalizujące i generyczne lambdy w C++14, przechwytywanie kopii
\texttt{*this} w C++17, a jawne listy parametrów szablonu lambdy
w C++20. Wykorzystamy semantykę przenoszenia z rozdziału~\ref{chap:move-semantics-cpp},
szablony wariadyczne z rozdziału~\ref{chap:variadic-templates-cpp}
i koncepcje z rozdziału~\ref{chap:concepts-cpp}.

\section{Trzy podstawowe postacie kodu wywoływalnego}
\label{sec:callables-forms}

\subsection{Funkcja i wskaźnik do funkcji}

Zwykła funkcja nie przechowuje kontekstu wywołania. Jej adres można
przypisać wskaźnikowi o odpowiednim typie:

\begin{lstlisting}
int twice(int x) { return 2 * x; }

int (*fp)(int) = &twice;
int a = fp(7);       // 14
int b = (*fp)(7);    // also 14
\end{lstlisting}

W deklaracji \texttt{int (*fp)(int)} nawiasy są istotne. Sam wskaźnik
przechowuje adres funkcji; nie ma miejsca na dodatkowy stan, taki jak
współczynnik mnożenia ustalany w czasie działania programu.

\subsection{Funktor: klasa z \texttt{operator()}}

\begin{lstlisting}
struct Multiply {
    int factor;

    int operator()(int x) const {
        return factor * x;
    }
};

Multiply times_three{3};
int result = times_three(7); // 21
\end{lstlisting}

\texttt{Multiply} jest klasą, a \texttt{times\_three} jej egzemplarzem.
Dzięki polu \texttt{factor} każdy egzemplarz może reprezentować inne
zachowanie. Operator wywołania może być przeciążany, szablonowy,
opatrzony kwalifikatorami \texttt{const} lub \texttt{\&\&} i mieć
złożoną implementację.

\subsection{Lambda jako wygodny zapis funktora}

\begin{lstlisting}
int factor = 3;
auto times_three = [factor](int x) {
    return factor * x;
};
int result = times_three(7); // 21
\end{lstlisting}

Kompilator tworzy odrębny, nienazwany typ klasy odpowiadający lambdzie.
Uproszczony model jej działania jest następujący:

\begin{lstlisting}
struct ClosureModel {
    int factor;

    int operator()(int x) const {
        return factor * x;
    }
};
\end{lstlisting}

Jest to \emph{model pojęciowy}, nie obiecana przez standard reprezentacja
pamięci. Każde wystąpienie wyrażenia lambda ma własny typ, nawet jeżeli
zapis kodu dwóch lambd wygląda identycznie.

\section{Budowa wyrażenia lambda}
\label{sec:lambda-syntax}

Najważniejsze elementy zapisu przedstawia schemat:

\begin{lstlisting}
[captures](parameters) mutable noexcept -> ResultType {
    // body
}
\end{lstlisting}

W C++20 można również jawnie określić parametry szablonu, koncepcje
oraz klauzulę \texttt{requires}. Typ wyniku zwykle jest wywnioskowany
z instrukcji \texttt{return}; podanie \texttt{-> ResultType} bywa
przydatne, gdy chcemy wymusić konkretny typ lub zwracać referencję.

\begin{lstlisting}
auto square = [](int x) { return x * x; };
auto half = [](double x) -> double { return x / 2; };
auto truth = [] { return true; }; // () may be omitted

int a = square(4);  // 16
double b = half(5); // 2.5
\end{lstlisting}

\paragraph{Najważniejsze zastrzeżenie.}
Lambda nie jest wskaźnikiem do funkcji. Zwykle tworzy obiekt z
\texttt{operator()}; jedynie pewne lambdy bez przechwyceń mają
konwersję na wskaźnik do funkcji o pasującej sygnaturze.

\section{Przechwytywanie zmiennych (\textit{captures})}

\begin{center}\small
\begin{tabularx}{\textwidth}{@{}p{0.30\textwidth}X@{}}
\toprule
Zapis lambdy & Znaczenie \\
\midrule
\code{[x]} & Przechwytuje wartość \code{x} przez kopię w obiekcie domknięcia. \\
\texttt{[\&x]} & Przechwytuje referencję do \code{x}; wywołanie wymaga, by \code{x} nadal żyło. \\
\code{[=]} & Domyślnie przechwytuje użyte zmienne lokalne przez wartość. \\
\texttt{[\&]} & Domyślnie przechwytuje użyte zmienne lokalne przez referencję. \\
\code{[p = std::move(p)]} & Przenosi wartość do składowej domknięcia przez inicjalizację przechwycenia. \\
\code{mutable} & Pozwala zmieniać kopie przechwycone przez wartość wewnątrz lambdy. \\
\bottomrule
\end{tabularx}
\end{center}
\label{sec:lambda-captures}

\subsection{Przechwycenie przez wartość}

\begin{lstlisting}
int threshold = 10;
auto above = [threshold](int x) {
    return x > threshold;
};
auto above_current = [&threshold](int x) {
    return x > threshold;
};
threshold = 100;

bool first = above(20);         // true: zapamietane 10
bool second = above_current(20); // false: aktualne 100
\end{lstlisting}

\texttt{[threshold]} inicjalizuje składową obiektu domknięcia wartością
\texttt{threshold} w chwili tworzenia lambdy. Późniejsza zmiana zmiennej
zewnętrznej nie modyfikuje tej kopii. Jeśli kopiowany obiekt ma kosztowne
kopiowanie, samo utworzenie lub skopiowanie lambdy może także mieć koszt.

\subsection{Przechwycenie przez referencję}

\begin{lstlisting}
int total = 0;
auto add = [&total](int value) {
    total += value;
};
add(4);
add(8); // total == 12
\end{lstlisting}

\texttt{[\&total]} uzyskuje dostęp do tej samej zmiennej, nie do jej
kopii. Lambda \textbf{nie wydłuża czasu życia} obiektu przechwyconego
przez referencję. Dlatego jest wygodna, gdy callback wykonywany jest
synchronicznie, ale potencjalnie niebezpieczna, jeśli jest zapamiętywany
albo wykonywany później.

\subsection{Przechwycenia domyślne i mieszane}

\begin{lstlisting}
int a = 2;
int b = 3;
int count = 0;

auto first  = [=]  { return a + b; }; // by copy
auto second = [&]  { ++count; };     // by reference
auto third  = [=, &count] {
    ++count;
    return a + b;
};
\end{lstlisting}

\texttt{[=]} i \texttt{[\&]} obejmują te użyte w treści zmienne lokalne,
które podlegają niejawnemu przechwyceniu; nie oznaczają automatycznego
skopiowania \emph{wszystkich} zmiennych z otoczenia. W większym kodzie
zwykle lepiej wypisać przechwycenia jawnie: dzięki temu widać zależności
i łatwiej ocenić bezpieczeństwo czasu życia.

\subsection{Przechwytywanie inicjalizujące (C++14)}

W nawiasach kwadratowych można zdefiniować nową nazwę i sposób jej
inicjalizacji. Dzięki temu lambdy mogą przechowywać wartości pochodne,
a także zasoby, których nie wolno kopiować:

\begin{lstlisting}
#include <memory>
#include <utility>

auto p = std::make_unique<int>(42);
auto read = [q = std::move(p)] {
    return *q;
};
// p == nullptr; ownership is now in the closure
int answer = read(); // 42
\end{lstlisting}

Taka lambda jest \emph{move-only}, ponieważ zawiera
\texttt{std::unique\_ptr}. Nie można jej skopiować. To istotne przy
wyborze miejsca przechowywania callbacków: w C++20
\texttt{std::function} wymaga kopiowalnego obiektu docelowego.

\subsection{Niebezpieczne przechwytywanie: wisząca referencja}

\begin{lstlisting}
// WRONG: do not call the returned lambda!
auto broken() {
    int local = 5;
    return [&local] { return local; };
}

// Correct: capture the value.
auto safe() {
    int local = 5;
    return [local] { return local; };
}
\end{lstlisting}

\texttt{broken()} zwraca obiekt z dostępem do zmiennej, która przestaje
istnieć wraz z wyjściem z funkcji. Późniejsze wywołanie prowadziłoby
do niezdefiniowanego zachowania (\textit{undefined behavior}).
Przechwycenie przez kopię rozwiązuje ten konkretny problem.

\section{Przechwytywanie \texttt{this} i \texttt{*this}}
\label{sec:lambda-this}

W metodzie klasy można korzystać z jej pól wewnątrz lambdy. Trzeba
jednak rozróżnić przechwycenie wskaźnika \texttt{this} od skopiowania
bieżącego obiektu.

\begin{lstlisting}
#include <string>

struct Label {
    std::string text;

    auto by_pointer() const {
        return [this] { return text; };
    }

    auto by_copy() const {
        return [*this] { return text; }; // C++17
    }
};
\end{lstlisting}

Lambda zwrócona przez \texttt{by\_pointer()} odwołuje się do oryginalnego
obiektu. Nie może być bezpiecznie wywoływana po jego zniszczeniu.
Lambda z \texttt{[*this]} przechowuje kopię obiektu (jeśli da się go
skopiować); dla składowych zarządzanych przez RAII jest to zwykle
bardziej przewidywalne, ale może być kosztowne.

\textbf{Uwaga:} zapis \texttt{[=]} użyty w metodzie nie kopiuje automatycznie
całego obiektu. Niejawne przechwytywanie \texttt{this} przez \texttt{[=]}
jest w C++20 przestarzałe (\textit{deprecated}); zapisujmy
\texttt{[this]} jawnie albo \texttt{[*this]}, gdy potrzebujemy kopii.

\section{Domyślnie stały operator wywołania i \texttt{mutable}}

Dla typowej lambdy bez słowa \texttt{mutable} operator wywołania ma
kwalifikator \texttt{const}. Oznacza to, że przechwyconych przez wartość
składowych nie wolno w nim zwyczajnie modyfikować:

\begin{lstlisting}
int counter = 0;

// ERROR: captured counter is not writable in this closure.
// auto bad = [counter] { ++counter; };

auto next = [counter]() mutable {
    return ++counter;
};

int a = next(); // 1
int b = next(); // 2
// outer counter is still 0
\end{lstlisting}

\texttt{mutable} modyfikuje zachowanie \emph{kopii wewnątrz lambdy},
nie powoduje przechwycenia oryginalnej zmiennej przez referencję.
Różne kopie obiektu domknięcia mogą przechowywać różny stan:

\begin{lstlisting}
auto a = [n = 0]() mutable { return ++n; };
auto b = a; // independent copy of state

a(); // 1
a(); // 2
b(); // 1
\end{lstlisting}

Przy przechwyceniu przez referencję do modyfikacji zewnętrznego
obiektu \texttt{mutable} nie jest potrzebne: stałość metody
\texttt{operator()} nie zmienia stałości obiektu osiąganego przez
referencję.

\section{Typ lambdy, kopiowanie i konwersja na wskaźnik}

\begin{lstlisting}
#include <type_traits>

inline constexpr auto f = [](int x) { return x + 1; };

auto a = [](int x) { return x; };
auto b = [](int x) { return x; };
static_assert(!std::is_same_v<decltype(a), decltype(b)>);

int (*pf)(int) = +[](int x) { return x + 1; };
\end{lstlisting}

Dwie osobno zapisane lambdy mają różne typy, nawet jeśli robią to samo.
W C++20 typ domknięcia lambdy bez przechwyceń jest domyślnie
konstruowalny; dostępność różnych operacji kopiowania/przypisywania
zależy od tego, czy lambda przechwytuje zmienne. Dla lambd z
przechwyceniami obowiązują też ograniczenia narzucone przez typy pól
(na przykład niekopiowalność \texttt{unique\_ptr}).

\texttt{+[]} jest znanym idiomem wymuszającym konwersję
nieprzechwytującej żadnej zmiennej lambdy na wskaźnik do funkcji.
Nie działa dla lambd przechwytujących:

\begin{lstlisting}
int offset = 10;
auto with_state = [offset](int x) { return x + offset; };
// int (*bad)(int) = with_state;  // ERROR
\end{lstlisting}

\section{Lambdy generyczne i własne parametry szablonu}
\label{sec:lambda-generic}

\subsection{Parametry \texttt{auto} (C++14)}

\begin{lstlisting}
#include <string>
auto print_twice = [](const auto& value) {
    return value + value;
};

auto x = print_twice(7);                 // 14
auto s = print_twice(std::string("ab")); // "abab"
\end{lstlisting}

Parametr \texttt{auto} sprawia, że \texttt{operator()} obiektu domknięcia
jest szablonem funkcji. Przy wywołaniu kompilator dedukuje odpowiednie
typy tak jak przy zwykłej funkcji szablonowej.

\subsection{Jawne parametry szablonu w lambdzie (C++20)}

\begin{lstlisting}
#include <concepts>

auto equal_types = []<class T>(const T& a, const T& b) {
    return a == b;
};

auto sum_integrals = []<std::integral... Ts>(Ts... values) {
    return (0 + ... + values);
};

bool yes = equal_types(3, 3);
int total = sum_integrals(1, 2, 3, 4); // 10
\end{lstlisting}

W przypadku \texttt{equal\_types} oba argumenty muszą dać ten sam
dedukowany typ \texttt{T}. W drugim przykładzie łączymy \emph{concepts}
z pakietem parametrów i wyrażeniem fold znanym z
rozdziałów~\ref{chap:variadic-templates-cpp} i~\ref{chap:concepts-cpp}.

\subsection{Uniwersalny wrapper z perfect forwarding}

\begin{lstlisting}
#include <functional>
#include <utility>

auto invoke_once = []<class F, class... Args>(
    F&& f, Args&&... args) -> decltype(auto) {
    return std::invoke(std::forward<F>(f),
                       std::forward<Args>(args)...);
};
\end{lstlisting}

\texttt{F\&\&} i \texttt{Args\&\&...} są tu referencjami przekazującymi,
a \texttt{decltype(auto)} zachowuje typ wyniku, w tym jego
referencyjność. W odróżnieniu od \texttt{std::move},
\texttt{std::forward} nie wymusza kategorii rvalue dla każdego argumentu.
Dzięki \texttt{std::invoke} wrapper działa nie tylko dla zwykłych
funkcji i lambd, lecz także dla wskaźników do składowych.

\section{Wymuszanie sygnatury: \texttt{noexcept}, \texttt{requires}, zwracanie referencji}

\begin{lstlisting}
#include <concepts>
#include <utility>

auto positive = [](int x) noexcept {
    return x > 0;
};

auto increment = []<std::integral T>(T x) -> T {
    return x + 1;
};

auto first = []<class T>(T& x) -> T& {
    return x;
};

int value = 3;
first(value) = 17; // modifies value
\end{lstlisting}

Ważne są dwie subtelności. Po pierwsze, \texttt{noexcept} jest obietnicą
programisty; jeśli wyjątek wydostanie się z funkcji oznaczonej
\texttt{noexcept}, wywołane zostanie \texttt{std::terminate}.
Po drugie, domyślna dedukcja wyniku lambdy deklarowanej bez
\texttt{decltype(auto)} lub jawnego \texttt{-> T\&} zwykle usuwa
referencję. Właściwe zwracanie referencji wymaga zatem świadomego
określenia typu wyniku i zapewnienia jego poprawnego czasu życia.

\section{\texttt{std::invoke}: jeden mechanizm wywołania}

\begin{center}\small
\begin{tabularx}{\textwidth}{@{}p{0.34\textwidth}X@{}}
\toprule
Narzędzie & Za co odpowiada \\
\midrule
\code{std::invoke(f, args...)} & Jednolicie wywołuje funkcję, lambdę lub wskaźnik do składowej. \\
\code{std::function<R(Args...)>} & Przechowuje kopiowalny obiekt wywoływalny o ustalonej sygnaturze przez wymazywanie typu. \\
\code{std::ref(x)} & Opakowuje referencję, gdy API normalnie kopiuje argument. \\
\code{std::bind_front(f, args...)} & Wiąże pierwsze argumenty funkcji i zwraca nowy obiekt wywoływalny. \\
\code{noexcept} przy lambdzie & Określa, że operator wywołania nie powinien pozwolić wyjątkowi opuścić funkcji. \\
\bottomrule
\end{tabularx}
\end{center}
\label{sec:std-invoke}

\texttt{std::invoke} z nagłówka \texttt{<functional>} wywołuje funkcję,
funktor lub lambdę, ale obsługuje także wskaźnik do metody i pola.

\begin{lstlisting}
#include <functional>

struct Box {
    int value = 10;
    int add(int x) const { return value + x; }
};

Box b;
int a = std::invoke([](int x) { return x * 2; }, 4);
int c = std::invoke(&Box::add, b, 7);      // 17
int& r = std::invoke(&Box::value, b);
r = 25;                                   // b.value == 25
\end{lstlisting}

To podstawowa cegiełka generycznych algorytmów. W nagłówku
\texttt{<type\_traits>} zdefiniowano m.in. \texttt{std::invoke\_result\_t},
\texttt{std::is\_invocable\_v} i \texttt{std::is\_nothrow\_invocable\_v},
które pozwalają sprawdzić w czasie kompilacji poprawność wywołania.

\begin{lstlisting}
#include <type_traits>

static_assert(std::is_invocable_v<decltype(positive), int>);
static_assert(!std::is_invocable_v<decltype(positive),
                                   std::string>);
static_assert(std::is_nothrow_invocable_v<decltype(positive),
                                          int>);
\end{lstlisting}

\section{\texttt{std::function}: jeden typ dla wielu zachowań}
\label{sec:std-function}

Lambdy o różnych typach nie mogą bezpośrednio zastępować się w jednej
zmiennej zadeklarowanej przez \texttt{auto}. Jeżeli chcemy przechowywać
\emph{dowolny} kopiowalny obiekt wywoływalny o zadanej sygnaturze,
możemy użyć \texttt{std::function<R(Args...)>}.

\begin{lstlisting}
#include <functional>
#include <string>

std::function<int(int)> operation =
    [](int x) { return 2 * x; };

int factor = 5;
operation = [factor](int x) { return factor * x; };
int result = operation(3); // 15

operation = nullptr;
if (!operation) {
    // no callable stored
}
\end{lstlisting}

\texttt{std::function} stosuje \emph{type erasure}: ukrywa rzeczywisty
typ obiektu docelowego za wspólnym interfejsem wywołania. Może
przechowywać funkcje, funktory i lambdy z przechwyceniami. Nie jest
odpowiednikiem \texttt{auto}: deklaracja \texttt{auto} zachowuje
konkretny typ i pozwala zwykle kompilatorowi łatwiej zoptymalizować
wywołanie.

\subsection{Pusty \texttt{std::function} i wyjątek}

\begin{lstlisting}
#include <functional>

std::function<void()> callback;
if (callback) {
    callback();
}
// Calling callback() while empty throws std::bad_function_call.
\end{lstlisting}

Sprawdzenie \texttt{if (callback)} jest zatem standardową częścią
obsługi opcjonalnych callbacków. Nie należy przechowywać
\texttt{std::function} bez potrzeby, jeśli konkretny typ wywoływalny
jest znany w czasie kompilacji.

\subsection{Koszty i ograniczenie kopiowalności}

Zwykle \texttt{std::function} wymaga pośredniego wywołania. Przechowywanie
większego domknięcia może też wymagać dynamicznej alokacji (zależy to
od implementacji i rozmiaru obiektu). To koszt elastyczności API, nie
argument przeciw używaniu \texttt{std::function} w każdym przypadku.

W C++20 przechowywany obiekt docelowy musi być kopiowalny, dlatego
następujący kod nie jest poprawny:

\begin{lstlisting}
#include <functional>
#include <memory>

// ERROR in C++20: the closure is not copy-constructible.
// std::function<int()> f =
//     [p = std::make_unique<int>(9)] { return *p; };
\end{lstlisting}

Od \textbf{C++23} biblioteka udostępnia
\texttt{std::move\_only\_function}, która służy do przechowywania
nienadających się do kopiowania obiektów wywoływalnych. W C++20
można natomiast przechowywać taki obiekt bezpośrednio jako \texttt{auto}
lub przekazać go do szablonu, który nie wymaga kopiowania.

\section{Co wybrać: \texttt{auto}, szablon, wskaźnik czy \texttt{std::function}?}
\label{sec:lambda-choice}

\begin{center}
\begin{tabular}{|p{3.0cm}|p{4.4cm}|p{5.0cm}|}
\hline
\textbf{Forma} & \textbf{Zalety} & \textbf{Ograniczenia} \\
\hline
\texttt{auto f = lambda} & Konkretny typ, bez konieczności wymazywania typu, naturalna obsługa stanu & Zmienna ma ustalony typ domknięcia \\
\hline
\texttt{template<class F>} & Działa z wieloma typami, umożliwia inline i forwarding & Definicja szablonu zwykle dostępna w nagłówku; możliwe powiększanie kodu \\
\hline
Wskaźnik do funkcji & Prosty nieposiadający stanu interfejs & Nie przyjmuje lambd z przechwyceniami \\
\hline
\texttt{std::function} & Jedna sygnatura, zmienne zachowanie w czasie wykonania & Type erasure; koszt pośredniego wywołania; wymóg kopiowalności celu w C++20 \\
\hline
\end{tabular}
\end{center}

\paragraph{Praktyczna reguła.}
Jeżeli callback wykorzystujemy natychmiast, wybierajmy zwykle szablon
lub \texttt{auto}. Jeżeli musimy przechowywać i później podmieniać różne
callbacki o jednej sygnaturze, \texttt{std::function} jest często
czytelniejsza niż ręczne modelowanie polimorfizmu.

\section{Lambdy ze standardowymi algorytmami}
\label{sec:lambda-algorithms}

\subsection{\texttt{find\_if}, \texttt{count\_if}, \texttt{transform}}

\begin{lstlisting}
#include <algorithm>
#include <vector>

std::vector<int> xs{1, 2, 3, 4, 5, 6};

int limit = 3;
auto it = std::find_if(xs.begin(), xs.end(),
                       [limit](int x) { return x > limit; });
// *it == 4

auto even_count = std::count_if(xs.begin(), xs.end(),
                                [](int x) { return x % 2 == 0; });
// even_count == 3

std::vector<int> squares(xs.size());
std::transform(xs.begin(), xs.end(), squares.begin(),
               [](int x) { return x * x; });
\end{lstlisting}

Dla \texttt{std::find\_if} trzeba przed dereferencją sprawdzić, czy
wynik jest różny od \texttt{xs.end()}. Dla \texttt{std::transform} zbiór
docelowy musi już zawierać wystarczającą liczbę elementów: samo
\texttt{reserve()} zwiększa \texttt{capacity()}, lecz nie \texttt{size()}.
Alternatywą jest pusty wektor i iterator \texttt{std::back\_inserter}.

\subsection{\texttt{std::ranges} i projekcje (C++20)}

\begin{lstlisting}
#include <algorithm>
#include <ranges>
#include <string>
#include <vector>

struct Person {
    std::string name;
    int age;
};

std::vector<Person> people{{"Ada", 31}, {"Bob", 18}, {"Eve", 25}};

std::ranges::sort(people, {}, &Person::age);
// sorted by age ascending

auto adults = people | std::views::filter(
    [](const Person& p) { return p.age >= 18; });
\end{lstlisting}

Trzeci argument \texttt{ranges::sort} to \emph{projekcja}:
\texttt{\&Person::age} pobiera klucz sortowania, a domyślny komparator
porównuje klucze. \texttt{views::filter} nie buduje automatycznie nowego
wektora; jest leniwym widokiem, który odwołuje się do danych źródłowych.
Przy widokach szczególnie uważnie trzeba kontrolować czas życia
kontenera i przechwyceń predykatu.

\section{\texttt{std::ref}, \texttt{std::bind\_front} i częściowe wiązanie argumentów}
\label{sec:bind-front}

\begin{lstlisting}
#include <functional>

int add(int x, int y) { return x + y; }

auto add_ten = std::bind_front(add, 10); // C++20
int x = add_ten(5);                      // 15
\end{lstlisting}

\texttt{std::bind\_front} przechowuje funkcję i ustalone argumenty,
a pozostałe przyjmuje podczas wywołania. W wielu zastosowaniach
równoważna lambda jest czytelniejsza:

\begin{lstlisting}
auto add_ten_lambda = [](int x) { return add(10, x); };
\end{lstlisting}

Domyślnie argumenty przechowywane przez \texttt{bind\_front} są
przechowywane przez wartość, z uwzględnieniem przenoszenia. Jeśli
potrzebujemy odwołania do oryginalnego obiektu, stosujemy
\texttt{std::ref} (a dla dostępu tylko do odczytu
\texttt{std::cref}), pamiętając o czasie jego życia:

\begin{lstlisting}
#include <functional>

struct Accumulator {
    int total = 0;
    void add(int n) { total += n; }
};

Accumulator acc;
auto update = std::bind_front(&Accumulator::add, std::ref(acc));
update(7); // acc.total == 7
\end{lstlisting}

Starsze \texttt{std::bind} jest dostępne, ale ma bardziej skomplikowane
reguły dotyczące argumentów i podstawiania. W nowym kodzie zwykle
preferujemy lambdy lub \texttt{std::bind\_front}.

\section{Praktyczny projekt: rejestr callbacków}
\label{sec:callbacks}

Rozważmy klasę powiadamiającą o zmianie wartości. Do przechowywania
procedury obsługi zdarzenia użyjemy \texttt{std::function}. Dzięki
temu można ją zastąpić inną nawet po utworzeniu obiektu.

\begin{lstlisting}
#include <functional>
#include <utility>

class Counter {
public:
    using Callback = std::function<void(int)>;

    void on_change(Callback callback) {
        callback_ = std::move(callback);
    }

    void increment() {
        ++value_;
        if (callback_) callback_(value_);
    }

    int value() const noexcept { return value_; }

private:
    int value_ = 0;
    Callback callback_;
};

Counter counter;
counter.on_change([](int value) {
    // report the new value
});
counter.increment();
\end{lstlisting}

Kod w pojedynczym wątku jest prosty i wystarczający dla wielu
interfejsów zdarzeniowych. W programie wielowątkowym równoczesna
zmiana \texttt{callback\_} oraz wykonywanie \texttt{increment()}
wymaga dodatkowej synchronizacji. \texttt{std::function} nie zapewnia
automatycznie ochrony wewnętrznego stanu przed wyścigami danych.

\subsection{Callback a czas życia właściciela}

Typowy przebieg to: rejestrujemy callback, właściciel danych zostaje
zniszczony, a ktoś później wywołuje callback. Jeśli lambda zachowała
samą referencję lub surowy wskaźnik, dostęp byłby wtedy błędny.
Słaby wskaźnik w poniższym przykładzie pozwala sprawdzić czas życia.

\begin{lstlisting}
#include <memory>

struct Worker : std::enable_shared_from_this<Worker> {
    int total = 0;

    auto callback() {
        std::weak_ptr<Worker> weak = shared_from_this();
        return [weak](int n) {
            if (auto self = weak.lock()) {
                self->total += n;
            }
        };
    }
};
\end{lstlisting}

Przykład zakłada, że obiekt \texttt{Worker} jest już zarządzany przez
\texttt{std::shared\_ptr} w chwili wywołania \texttt{callback()}.
\texttt{weak\_ptr} pozwala sprawdzić, czy obiekt wciąż istnieje,
bez przedłużania jego życia przez przechwycenie silnego wskaźnika.
To ważne również dla unikania cykli własności, np. gdy sam obiekt
przechowuje callback odnoszący się do niego.

\section{Rekurencyjne lambdy bez \texttt{std::function}}

Lambda zwykle nie zna bezpośrednio własnej nazwy podczas definiowania.
Można jednak przekazywać obiekt wywoływalny jako argument:

\begin{lstlisting}
auto factorial = [](auto&& self, unsigned n) -> unsigned {
    return n <= 1 ? 1 : n * self(self, n - 1);
};

unsigned x = factorial(factorial, 5); // 120
\end{lstlisting}

To odmiana wzorca \emph{self-parameter}. Dla realnych zastosowań
rekurencyjnych warto ponadto zwrócić uwagę na przepełnienie liczb
całkowitych i maksymalną głębokość stosu. W C++20 metoda ta pozwala
uniknąć pośredniej alokacji i wymazywania typu, których mogłaby wymagać
rekurencja oparta na \texttt{std::function}.

\section{Lambdy w metaprogramowaniu: \texttt{std::visit} i \texttt{std::apply}}
\label{sec:lambda-visit-apply}

Dwa przydatne zastosowania łączą lambdy z mechanizmami z
rozdziałów~\ref{chap:cpp14-type-erasure} i~\ref{chap:variadic-templates-cpp}.
\texttt{std::visit} wybiera wywołanie na podstawie aktywnego typu
\texttt{std::variant}, natomiast \texttt{std::apply} rozwija argumenty
z krotki lub podobnej struktury.

\subsection{Zestaw przeciążeń oparty na dziedziczeniu}

\begin{lstlisting}
#include <string>
#include <variant>

template <class... Fs>
struct overloaded : Fs... {
    using Fs::operator()...;
};

std::variant<int, std::string> value = std::string{"Hello"};

auto stringify = overloaded{
    [](int x) { return std::to_string(x); },
    [](const std::string& x) { return x; }
};

std::string result = std::visit(stringify, value); // "Hello"
\end{lstlisting}

Każda lambda dostarcza jeden \texttt{operator()}. Dziedziczenie
wariadyczne i deklaracja \texttt{using Fs::operator()...} tworzą zbiór
przeciążeń, spośród których kompilator wybiera odpowiednią funkcję.
W tym przykładzie wszystkie warianty wywołania zwracają ten sam typ,
\texttt{std::string}, czego wymaga tu pojedyncze wywołanie
\texttt{std::visit}. Konstrukcja jest przydatna przy przetwarzaniu
stanów programu i heterogenicznych danych.

\subsection{Rozwijanie krotki przez \texttt{std::apply}}

\begin{lstlisting}
#include <tuple>

std::tuple<int, int, int> data{2, 3, 5};
int sum = std::apply(
    [](auto... values) { return (0 + ... + values); },
    data);
// sum == 10
\end{lstlisting}

Tutaj \texttt{std::apply} bierze elementy krotki i wywołuje lambdę
jak zwykłą funkcję trójargumentową. Sama lambda wykorzystuje
\emph{fold expression}, a więc nie trzeba pisać rekurencyjnego kodu.

\section{Pakiety przechwyceń i obliczenia odroczone (C++20)}

C++20 pozwala stosować \emph{pack expansion} również w przechwytywaniu
inicjalizującym. Możemy w ten sposób zbudować jeden obiekt domknięcia,
który przechowuje cały pakiet wartości do późniejszego przetworzenia:

\begin{lstlisting}
#include <concepts>

template <std::integral... Ts>
auto delayed_sum(Ts... xs) {
    return [...stored = xs] {
        return (0 + ... + stored);
    };
}

auto job = delayed_sum(2, 3, 5);
int result = job(); // 10
\end{lstlisting}

Lambdę można zwrócić z funkcji, ponieważ wartości są przechowywane
wewnątrz domknięcia. Przechwycenie przez referencję miałoby inne
wymagania dotyczące czasu życia. W rzeczywistych fabrykach callbacków
pakiet może zawierać obiekty niekopiowalne; wtedy należy zadbać
o przekazanie własności (np. przez \texttt{std::move} i
\texttt{std::forward}) oraz o wymagania interfejsu, który callback
przechowuje.

\section{Lambdy w obliczeniach stałych}

Jeśli ciało lambdy dopuszcza obliczanie stałe, jej wywołanie może
być częścią wyrażenia stałego. Od C++17 odpowiedni operator wywołania
jest domyślnie \texttt{constexpr}, gdy spełnia wymagania takiej funkcji:

\begin{lstlisting}
constexpr auto cube = [](int x) { return x * x * x; };
static_assert(cube(3) == 27);

// C++20: each call must be evaluated at compile time.
auto fourth_power = [](int x) consteval {
    return x * x * x * x;
};
static_assert(fourth_power(2) == 16);
\end{lstlisting}

\texttt{consteval} jest mocniejszą deklaracją niż \texttt{constexpr}:
wymaga, aby właściwe wywołania funkcji natychmiastowej spełniały
reguły obliczania stałego. Nie należy traktować tego jako gwarancji,
że każda zwykła lambda zostanie wykonana podczas kompilacji.

\section{Koncepcje dla obiektów wywoływalnych}
\label{sec:callable-concepts}

W C++20 można ograniczać typy wywoływalne m.in. następującymi
koncepcjami:
\begin{itemize}
\item \texttt{std::invocable} -- poprawność wywołania z danymi argumentami;
\item \texttt{std::regular\_invocable} -- dodatkowe wymagania semantyczne;
\item \texttt{std::predicate} -- poprawnie wywoływalny predykat z
wynikiem dającym się traktować jako wartość logiczną.
\end{itemize}
Koncepcje nie dowodzą wszystkich własności semantycznych. Przykładowo
kompilator nie zweryfikuje matematycznie stabilności wyników predykatu.

\begin{lstlisting}
#include <concepts>
#include <functional>
#include <utility>

template <class F, class T>
    requires std::predicate<F&, const T&>
bool test(const T& value, F&& predicate) {
    return std::invoke(std::forward<F>(predicate), value);
}

bool ok = test(8, [](int x) { return x % 2 == 0; });
\end{lstlisting}

\textbf{Istotne rozróżnienie:} koncepcja weryfikuje przede wszystkim
własności typów i poprawność wyrażeń. Nie dowodzi matematycznie
własności predykatu, np. że comparator tworzy ścisły słaby porządek.
Programista nadal odpowiada za kontrakt semantyczny.

\section{Typowe błędy i trudniejsze przypadki}
\label{sec:lambda-pitfalls}

\subsection{Referencja przechwycona do operacji asynchronicznej}

\begin{lstlisting}
#include <future>

std::future<int> start() {
    int x = 7;
    // WRONG: the task could outlive x.
    // return std::async(std::launch::async, [&x] { return x; });
    return std::async(std::launch::async, [x] { return x; });
}
\end{lstlisting}

W zadaniach asynchronicznych preferuj przechwycenia przez wartość
lub świadomie zarządzaną własność, dopóki nie mamy ścisłej gwarancji,
że obiekt referowany przetrwa wykonanie zadania. Sama biblioteka
\texttt{std::async} nie wydłuża życia lokalnych zmiennych funkcji
uruchamiającej.

\subsection{Predykat sortowania nie może naruszać porządku}

\begin{lstlisting}
// WRONG: <= is not a strict ordering relation.
// std::sort(xs.begin(), xs.end(),
//           [](int a, int b) { return a <= b; });

std::sort(xs.begin(), xs.end(),
          [](int a, int b) { return a < b; });
\end{lstlisting}

Komparator do \texttt{sort} musi definiować \emph{ścisły słaby
porządek}. Przykładowo \texttt{comp(x,x)} ma dawać \texttt{false}.
Naruszenie wymagań algorytmu może prowadzić do niezdefiniowanego
zachowania, mimo że typy argumentów się zgadzają.

\subsection{Pomylenie kopii z oryginalnym stanem}

\begin{lstlisting}
int sum = 0;
std::for_each(xs.begin(), xs.end(), [sum](int x) mutable {
    sum += x; // modifies only the closure's copy
});
// outer sum remains 0

std::for_each(xs.begin(), xs.end(), [&sum](int x) {
    sum += x;
});
\end{lstlisting}

Nie należy zakładać, że \texttt{mutable} zmienia zewnętrzną
zmienną. Należy też pamiętać, że algorytmy
mogą kopiować obiekt funkcyjny; nie opierajmy semantyki algorytmu
na nieudokumentowanych szczegółach kopiowania predykatu.

\subsection{Przechwycenie wskaźnika nie przedłuża życia obiektu}

\begin{lstlisting}
int* p = /* pointer to some object */ nullptr;
auto f = [p] {
    if (p) return *p;
    return 0;
};
\end{lstlisting}

Wartość wskaźnika może zostać skopiowana do lambdy, ale wskazywany
obiekt nadal musi istnieć w chwili dereferencji. Przechwycenie
\texttt{shared\_ptr} przenosi własność współdzieloną, a
\texttt{weak\_ptr} pozwala ją warunkowo odzyskać. To nie jest to
samo co skopiowanie wskaźnika surowego.

\section{Zestawienie: kiedy stosować którą konstrukcję?}

\begin{center}
\begin{tabular}{|p{5.7cm}|p{6.6cm}|}
\hline
\textbf{Potrzeba} & \textbf{Zalecany wybór} \\
\hline
Krótki predykat do algorytmu & Lambda z jawnymi przechwyceniami \\
\hline
Wywołanie bez stanu w API C & Zwykła funkcja albo zgodna lambda bez przechwyceń jako wskaźnik do funkcji \\
\hline
Wydajny callback o typie znanym w czasie kompilacji & Parametr szablonowy \texttt{F}, ewentualnie \texttt{std::invocable} \\
\hline
Callback z możliwością późniejszej podmiany & \texttt{std::function} (C++20, kopiowalny cel) \\
\hline
Przechowanie \texttt{unique\_ptr} wewnątrz lambdy & Init-capture i przeniesienie; nie przechowuj takiej lambdy w \texttt{std::function} w C++20 \\
\hline
Callback żyjący dłużej niż zmienne lokalne & Kopiuj potrzebny stan lub zarządzaj życiem przez RAII i smart pointery \\
\hline
Uniwersalne wywołanie dowolnego callable & \texttt{std::invoke} i \texttt{std::forward} \\
\hline
Częściowe podstawienie kilku argumentów & Lambda lub \texttt{std::bind\_front} \\
\hline
Bezpieczny predykat dla szablonu & \texttt{std::predicate} plus kontrola jego semantyki \\
\hline
\end{tabular}
\end{center}

\paragraph{Pytania kontrolne.}
\begin{enumerate}
\item Dlaczego dwie lambdy o identycznym kodzie mają różne typy?
\item Jaka jest różnica między \texttt{[x]}, \texttt{[\&x]} i \texttt{[x]() mutable}?
\item Kiedy bezpiecznie zwrócić lambdę przechwytującą zmienną lokalną?
\item Dlaczego lambda posiadająca \texttt{unique\_ptr} nie może być celem \texttt{std::function} w C++20?
\item Jak \texttt{std::invoke} wywołuje wskaźnik do metody klasy?
\item Czym różnią się \texttt{[this]} i \texttt{[*this]}?
\item Co może pójść źle, gdy \texttt{std::views::filter} korzysta z lokalnej zmiennej przez referencję?
\item Kiedy interfejs przyjmujący szablon \texttt{F} jest lepszy od \texttt{std::function}?
\item Dlaczego komparator \texttt{a <= b} nie jest poprawny dla \texttt{std::sort}?
\item W jakiej sytuacji przechwycenie \texttt{weak\_ptr} jest lepsze od \texttt{shared\_ptr}?
\end{enumerate}

\paragraph{Dalsza lektura.}
W dokumentacji \emph{cppreference} warto przeczytać hasła
\emph{Lambda expressions}, \emph{Function objects} i \emph{Ranges}.
Dokładne opisy narzędzi znajdują się również pod hasłami:
\begin{itemize}
\item \texttt{std::invoke} oraz \texttt{std::function};
\item \texttt{std::bind\_front} oraz \texttt{std::predicate}.
\end{itemize}


\chapter[Lambdy i obiekty funkcyjne]{Wyrażenia lambda, obiekty funkcyjne i abstrakcje wywoływalne}
\label{chap:lambdas-cpp}

\section{Po co istnieją obiekty funkcyjne?}

W C++ funkcję można reprezentować nie tylko przez jej nazwę lub wskaźnik.
Można również utworzyć \emph{obiekt wywoływalny} (\textit{callable object}):
obiekt, który daje się zastosować w wyrażeniu przypominającym wywołanie
funkcji. Pozwala to przekazywać zachowanie jako parametr algorytmu,
zapamiętywać stan pomiędzy wywołaniami oraz tworzyć elastyczne API
(np. sortowanie z własnym porządkiem, filtrowanie, obsługa zdarzeń).

\begin{lstlisting}
#include <algorithm>
#include <vector>

std::vector<int> numbers{5, 2, 8, 1};
std::sort(numbers.begin(), numbers.end(),
          [](int a, int b) { return a > b; });
// numbers == {8, 5, 2, 1}
\end{lstlisting}

Trzeci argument \texttt{std::sort} to reguła porównywania. Zamiast pisać
osobną funkcję o globalnej nazwie, definiujemy ją dokładnie w miejscu
użycia. Wyrażenie \texttt{[](int a, int b) \{ return a > b; \}}
jest \emph{lambdą}: tworzy obiekt unikatowego typu domknięcia
(\textit{closure type}).

\paragraph{Zakres i wymagania.}
Przyjmujemy \textbf{C++20}. Lambdy wprowadzono w C++11, przechwytywanie
inicjalizujące i generyczne lambdy w C++14, przechwytywanie kopii
\texttt{*this} w C++17, a jawne listy parametrów szablonu lambdy
w C++20. Wykorzystamy semantykę przenoszenia z rozdziału~\ref{chap:move-semantics-cpp},
szablony wariadyczne z rozdziału~\ref{chap:variadic-templates-cpp}
i koncepcje z rozdziału~\ref{chap:concepts-cpp}.

\section{Trzy podstawowe postacie kodu wywoływalnego}
\label{sec:callables-forms}

\subsection{Funkcja i wskaźnik do funkcji}

Zwykła funkcja nie przechowuje kontekstu wywołania. Jej adres można
przypisać wskaźnikowi o odpowiednim typie:

\begin{lstlisting}
int twice(int x) { return 2 * x; }

int (*fp)(int) = &twice;
int a = fp(7);       // 14
int b = (*fp)(7);    // also 14
\end{lstlisting}

W deklaracji \texttt{int (*fp)(int)} nawiasy są istotne. Sam wskaźnik
przechowuje adres funkcji; nie ma miejsca na dodatkowy stan, taki jak
współczynnik mnożenia ustalany w czasie działania programu.

\subsection{Funktor: klasa z \texttt{operator()}}

\begin{lstlisting}
struct Multiply {
    int factor;

    int operator()(int x) const {
        return factor * x;
    }
};

Multiply times_three{3};
int result = times_three(7); // 21
\end{lstlisting}

\texttt{Multiply} jest klasą, a \texttt{times\_three} jej egzemplarzem.
Dzięki polu \texttt{factor} każdy egzemplarz może reprezentować inne
zachowanie. Operator wywołania może być przeciążany, szablonowy,
opatrzony kwalifikatorami \texttt{const} lub \texttt{\&\&} i mieć
złożoną implementację.

\subsection{Lambda jako wygodny zapis funktora}

\begin{lstlisting}
int factor = 3;
auto times_three = [factor](int x) {
    return factor * x;
};
int result = times_three(7); // 21
\end{lstlisting}

Kompilator tworzy odrębny, nienazwany typ klasy odpowiadający lambdzie.
Uproszczony model jej działania jest następujący:

\begin{lstlisting}
struct ClosureModel {
    int factor;

    int operator()(int x) const {
        return factor * x;
    }
};
\end{lstlisting}

Jest to \emph{model pojęciowy}, nie obiecana przez standard reprezentacja
pamięci. Każde wystąpienie wyrażenia lambda ma własny typ, nawet jeżeli
zapis kodu dwóch lambd wygląda identycznie.

\section{Budowa wyrażenia lambda}
\label{sec:lambda-syntax}

Najważniejsze elementy zapisu przedstawia schemat:

\begin{lstlisting}
[captures](parameters) mutable noexcept -> ResultType {
    // body
}
\end{lstlisting}

W C++20 można również jawnie określić parametry szablonu, koncepcje
oraz klauzulę \texttt{requires}. Typ wyniku zwykle jest wywnioskowany
z instrukcji \texttt{return}; podanie \texttt{-> ResultType} bywa
przydatne, gdy chcemy wymusić konkretny typ lub zwracać referencję.

\begin{lstlisting}
auto square = [](int x) { return x * x; };
auto half = [](double x) -> double { return x / 2; };
auto truth = [] { return true; }; // () may be omitted

int a = square(4);  // 16
double b = half(5); // 2.5
\end{lstlisting}

\paragraph{Najważniejsze zastrzeżenie.}
Lambda nie jest wskaźnikiem do funkcji. Zwykle tworzy obiekt z
\texttt{operator()}; jedynie pewne lambdy bez przechwyceń mają
konwersję na wskaźnik do funkcji o pasującej sygnaturze.

\section{Przechwytywanie zmiennych (\textit{captures})}

\begin{center}\small
\begin{tabularx}{\textwidth}{@{}p{0.30\textwidth}X@{}}
\toprule
Zapis lambdy & Znaczenie \\
\midrule
\code{[x]} & Przechwytuje wartość \code{x} przez kopię w obiekcie domknięcia. \\
\texttt{[\&x]} & Przechwytuje referencję do \code{x}; wywołanie wymaga, by \code{x} nadal żyło. \\
\code{[=]} & Domyślnie przechwytuje użyte zmienne lokalne przez wartość. \\
\texttt{[\&]} & Domyślnie przechwytuje użyte zmienne lokalne przez referencję. \\
\code{[p = std::move(p)]} & Przenosi wartość do składowej domknięcia przez inicjalizację przechwycenia. \\
\code{mutable} & Pozwala zmieniać kopie przechwycone przez wartość wewnątrz lambdy. \\
\bottomrule
\end{tabularx}
\end{center}
\label{sec:lambda-captures}

\subsection{Przechwycenie przez wartość}

\begin{lstlisting}
int threshold = 10;
auto above = [threshold](int x) {
    return x > threshold;
};
auto above_current = [&threshold](int x) {
    return x > threshold;
};
threshold = 100;

bool first = above(20);         // true: zapamietane 10
bool second = above_current(20); // false: aktualne 100
\end{lstlisting}

\texttt{[threshold]} inicjalizuje składową obiektu domknięcia wartością
\texttt{threshold} w chwili tworzenia lambdy. Późniejsza zmiana zmiennej
zewnętrznej nie modyfikuje tej kopii. Jeśli kopiowany obiekt ma kosztowne
kopiowanie, samo utworzenie lub skopiowanie lambdy może także mieć koszt.

\subsection{Przechwycenie przez referencję}

\begin{lstlisting}
int total = 0;
auto add = [&total](int value) {
    total += value;
};
add(4);
add(8); // total == 12
\end{lstlisting}

\texttt{[\&total]} uzyskuje dostęp do tej samej zmiennej, nie do jej
kopii. Lambda \textbf{nie wydłuża czasu życia} obiektu przechwyconego
przez referencję. Dlatego jest wygodna, gdy callback wykonywany jest
synchronicznie, ale potencjalnie niebezpieczna, jeśli jest zapamiętywany
albo wykonywany później.

\subsection{Przechwycenia domyślne i mieszane}

\begin{lstlisting}
int a = 2;
int b = 3;
int count = 0;

auto first  = [=]  { return a + b; }; // by copy
auto second = [&]  { ++count; };     // by reference
auto third  = [=, &count] {
    ++count;
    return a + b;
};
\end{lstlisting}

\texttt{[=]} i \texttt{[\&]} obejmują te użyte w treści zmienne lokalne,
które podlegają niejawnemu przechwyceniu; nie oznaczają automatycznego
skopiowania \emph{wszystkich} zmiennych z otoczenia. W większym kodzie
zwykle lepiej wypisać przechwycenia jawnie: dzięki temu widać zależności
i łatwiej ocenić bezpieczeństwo czasu życia.

\subsection{Przechwytywanie inicjalizujące (C++14)}

W nawiasach kwadratowych można zdefiniować nową nazwę i sposób jej
inicjalizacji. Dzięki temu lambdy mogą przechowywać wartości pochodne,
a także zasoby, których nie wolno kopiować:

\begin{lstlisting}
#include <memory>
#include <utility>

auto p = std::make_unique<int>(42);
auto read = [q = std::move(p)] {
    return *q;
};
// p == nullptr; ownership is now in the closure
int answer = read(); // 42
\end{lstlisting}

Taka lambda jest \emph{move-only}, ponieważ zawiera
\texttt{std::unique\_ptr}. Nie można jej skopiować. To istotne przy
wyborze miejsca przechowywania callbacków: w C++20
\texttt{std::function} wymaga kopiowalnego obiektu docelowego.

\subsection{Niebezpieczne przechwytywanie: wisząca referencja}

\begin{lstlisting}
// WRONG: do not call the returned lambda!
auto broken() {
    int local = 5;
    return [&local] { return local; };
}

// Correct: capture the value.
auto safe() {
    int local = 5;
    return [local] { return local; };
}
\end{lstlisting}

\texttt{broken()} zwraca obiekt z dostępem do zmiennej, która przestaje
istnieć wraz z wyjściem z funkcji. Późniejsze wywołanie prowadziłoby
do niezdefiniowanego zachowania (\textit{undefined behavior}).
Przechwycenie przez kopię rozwiązuje ten konkretny problem.

\section{Przechwytywanie \texttt{this} i \texttt{*this}}
\label{sec:lambda-this}

W metodzie klasy można korzystać z jej pól wewnątrz lambdy. Trzeba
jednak rozróżnić przechwycenie wskaźnika \texttt{this} od skopiowania
bieżącego obiektu.

\begin{lstlisting}
#include <string>

struct Label {
    std::string text;

    auto by_pointer() const {
        return [this] { return text; };
    }

    auto by_copy() const {
        return [*this] { return text; }; // C++17
    }
};
\end{lstlisting}

Lambda zwrócona przez \texttt{by\_pointer()} odwołuje się do oryginalnego
obiektu. Nie może być bezpiecznie wywoływana po jego zniszczeniu.
Lambda z \texttt{[*this]} przechowuje kopię obiektu (jeśli da się go
skopiować); dla składowych zarządzanych przez RAII jest to zwykle
bardziej przewidywalne, ale może być kosztowne.

\textbf{Uwaga:} zapis \texttt{[=]} użyty w metodzie nie kopiuje automatycznie
całego obiektu. Niejawne przechwytywanie \texttt{this} przez \texttt{[=]}
jest w C++20 przestarzałe (\textit{deprecated}); zapisujmy
\texttt{[this]} jawnie albo \texttt{[*this]}, gdy potrzebujemy kopii.

\section{Domyślnie stały operator wywołania i \texttt{mutable}}

Dla typowej lambdy bez słowa \texttt{mutable} operator wywołania ma
kwalifikator \texttt{const}. Oznacza to, że przechwyconych przez wartość
składowych nie wolno w nim zwyczajnie modyfikować:

\begin{lstlisting}
int counter = 0;

// ERROR: captured counter is not writable in this closure.
// auto bad = [counter] { ++counter; };

auto next = [counter]() mutable {
    return ++counter;
};

int a = next(); // 1
int b = next(); // 2
// outer counter is still 0
\end{lstlisting}

\texttt{mutable} modyfikuje zachowanie \emph{kopii wewnątrz lambdy},
nie powoduje przechwycenia oryginalnej zmiennej przez referencję.
Różne kopie obiektu domknięcia mogą przechowywać różny stan:

\begin{lstlisting}
auto a = [n = 0]() mutable { return ++n; };
auto b = a; // independent copy of state

a(); // 1
a(); // 2
b(); // 1
\end{lstlisting}

Przy przechwyceniu przez referencję do modyfikacji zewnętrznego
obiektu \texttt{mutable} nie jest potrzebne: stałość metody
\texttt{operator()} nie zmienia stałości obiektu osiąganego przez
referencję.

\section{Typ lambdy, kopiowanie i konwersja na wskaźnik}

\begin{lstlisting}
#include <type_traits>

inline constexpr auto f = [](int x) { return x + 1; };

auto a = [](int x) { return x; };
auto b = [](int x) { return x; };
static_assert(!std::is_same_v<decltype(a), decltype(b)>);

int (*pf)(int) = +[](int x) { return x + 1; };
\end{lstlisting}

Dwie osobno zapisane lambdy mają różne typy, nawet jeśli robią to samo.
W C++20 typ domknięcia lambdy bez przechwyceń jest domyślnie
konstruowalny; dostępność różnych operacji kopiowania/przypisywania
zależy od tego, czy lambda przechwytuje zmienne. Dla lambd z
przechwyceniami obowiązują też ograniczenia narzucone przez typy pól
(na przykład niekopiowalność \texttt{unique\_ptr}).

\texttt{+[]} jest znanym idiomem wymuszającym konwersję
nieprzechwytującej żadnej zmiennej lambdy na wskaźnik do funkcji.
Nie działa dla lambd przechwytujących:

\begin{lstlisting}
int offset = 10;
auto with_state = [offset](int x) { return x + offset; };
// int (*bad)(int) = with_state;  // ERROR
\end{lstlisting}

\section{Lambdy generyczne i własne parametry szablonu}
\label{sec:lambda-generic}

\subsection{Parametry \texttt{auto} (C++14)}

\begin{lstlisting}
#include <string>
auto print_twice = [](const auto& value) {
    return value + value;
};

auto x = print_twice(7);                 // 14
auto s = print_twice(std::string("ab")); // "abab"
\end{lstlisting}

Parametr \texttt{auto} sprawia, że \texttt{operator()} obiektu domknięcia
jest szablonem funkcji. Przy wywołaniu kompilator dedukuje odpowiednie
typy tak jak przy zwykłej funkcji szablonowej.

\subsection{Jawne parametry szablonu w lambdzie (C++20)}

\begin{lstlisting}
#include <concepts>

auto equal_types = []<class T>(const T& a, const T& b) {
    return a == b;
};

auto sum_integrals = []<std::integral... Ts>(Ts... values) {
    return (0 + ... + values);
};

bool yes = equal_types(3, 3);
int total = sum_integrals(1, 2, 3, 4); // 10
\end{lstlisting}

W przypadku \texttt{equal\_types} oba argumenty muszą dać ten sam
dedukowany typ \texttt{T}. W drugim przykładzie łączymy \emph{concepts}
z pakietem parametrów i wyrażeniem fold znanym z
rozdziałów~\ref{chap:variadic-templates-cpp} i~\ref{chap:concepts-cpp}.

\subsection{Uniwersalny wrapper z perfect forwarding}

\begin{lstlisting}
#include <functional>
#include <utility>

auto invoke_once = []<class F, class... Args>(
    F&& f, Args&&... args) -> decltype(auto) {
    return std::invoke(std::forward<F>(f),
                       std::forward<Args>(args)...);
};
\end{lstlisting}

\texttt{F\&\&} i \texttt{Args\&\&...} są tu referencjami przekazującymi,
a \texttt{decltype(auto)} zachowuje typ wyniku, w tym jego
referencyjność. W odróżnieniu od \texttt{std::move},
\texttt{std::forward} nie wymusza kategorii rvalue dla każdego argumentu.
Dzięki \texttt{std::invoke} wrapper działa nie tylko dla zwykłych
funkcji i lambd, lecz także dla wskaźników do składowych.

\section{Wymuszanie sygnatury: \texttt{noexcept}, \texttt{requires}, zwracanie referencji}

\begin{lstlisting}
#include <concepts>
#include <utility>

auto positive = [](int x) noexcept {
    return x > 0;
};

auto increment = []<std::integral T>(T x) -> T {
    return x + 1;
};

auto first = []<class T>(T& x) -> T& {
    return x;
};

int value = 3;
first(value) = 17; // modifies value
\end{lstlisting}

Ważne są dwie subtelności. Po pierwsze, \texttt{noexcept} jest obietnicą
programisty; jeśli wyjątek wydostanie się z funkcji oznaczonej
\texttt{noexcept}, wywołane zostanie \texttt{std::terminate}.
Po drugie, domyślna dedukcja wyniku lambdy deklarowanej bez
\texttt{decltype(auto)} lub jawnego \texttt{-> T\&} zwykle usuwa
referencję. Właściwe zwracanie referencji wymaga zatem świadomego
określenia typu wyniku i zapewnienia jego poprawnego czasu życia.

\section{\texttt{std::invoke}: jeden mechanizm wywołania}

\begin{center}\small
\begin{tabularx}{\textwidth}{@{}p{0.34\textwidth}X@{}}
\toprule
Narzędzie & Za co odpowiada \\
\midrule
\code{std::invoke(f, args...)} & Jednolicie wywołuje funkcję, lambdę lub wskaźnik do składowej. \\
\code{std::function<R(Args...)>} & Przechowuje kopiowalny obiekt wywoływalny o ustalonej sygnaturze przez wymazywanie typu. \\
\code{std::ref(x)} & Opakowuje referencję, gdy API normalnie kopiuje argument. \\
\code{std::bind_front(f, args...)} & Wiąże pierwsze argumenty funkcji i zwraca nowy obiekt wywoływalny. \\
\code{noexcept} przy lambdzie & Określa, że operator wywołania nie powinien pozwolić wyjątkowi opuścić funkcji. \\
\bottomrule
\end{tabularx}
\end{center}
\label{sec:std-invoke}

\texttt{std::invoke} z nagłówka \texttt{<functional>} wywołuje funkcję,
funktor lub lambdę, ale obsługuje także wskaźnik do metody i pola.

\begin{lstlisting}
#include <functional>

struct Box {
    int value = 10;
    int add(int x) const { return value + x; }
};

Box b;
int a = std::invoke([](int x) { return x * 2; }, 4);
int c = std::invoke(&Box::add, b, 7);      // 17
int& r = std::invoke(&Box::value, b);
r = 25;                                   // b.value == 25
\end{lstlisting}

To podstawowa cegiełka generycznych algorytmów. W nagłówku
\texttt{<type\_traits>} zdefiniowano m.in. \texttt{std::invoke\_result\_t},
\texttt{std::is\_invocable\_v} i \texttt{std::is\_nothrow\_invocable\_v},
które pozwalają sprawdzić w czasie kompilacji poprawność wywołania.

\begin{lstlisting}
#include <type_traits>

static_assert(std::is_invocable_v<decltype(positive), int>);
static_assert(!std::is_invocable_v<decltype(positive),
                                   std::string>);
static_assert(std::is_nothrow_invocable_v<decltype(positive),
                                          int>);
\end{lstlisting}

\section{\texttt{std::function}: jeden typ dla wielu zachowań}
\label{sec:std-function}

Lambdy o różnych typach nie mogą bezpośrednio zastępować się w jednej
zmiennej zadeklarowanej przez \texttt{auto}. Jeżeli chcemy przechowywać
\emph{dowolny} kopiowalny obiekt wywoływalny o zadanej sygnaturze,
możemy użyć \texttt{std::function<R(Args...)>}.

\begin{lstlisting}
#include <functional>
#include <string>

std::function<int(int)> operation =
    [](int x) { return 2 * x; };

int factor = 5;
operation = [factor](int x) { return factor * x; };
int result = operation(3); // 15

operation = nullptr;
if (!operation) {
    // no callable stored
}
\end{lstlisting}

\texttt{std::function} stosuje \emph{type erasure}: ukrywa rzeczywisty
typ obiektu docelowego za wspólnym interfejsem wywołania. Może
przechowywać funkcje, funktory i lambdy z przechwyceniami. Nie jest
odpowiednikiem \texttt{auto}: deklaracja \texttt{auto} zachowuje
konkretny typ i pozwala zwykle kompilatorowi łatwiej zoptymalizować
wywołanie.

\subsection{Pusty \texttt{std::function} i wyjątek}

\begin{lstlisting}
#include <functional>

std::function<void()> callback;
if (callback) {
    callback();
}
// Calling callback() while empty throws std::bad_function_call.
\end{lstlisting}

Sprawdzenie \texttt{if (callback)} jest zatem standardową częścią
obsługi opcjonalnych callbacków. Nie należy przechowywać
\texttt{std::function} bez potrzeby, jeśli konkretny typ wywoływalny
jest znany w czasie kompilacji.

\subsection{Koszty i ograniczenie kopiowalności}

Zwykle \texttt{std::function} wymaga pośredniego wywołania. Przechowywanie
większego domknięcia może też wymagać dynamicznej alokacji (zależy to
od implementacji i rozmiaru obiektu). To koszt elastyczności API, nie
argument przeciw używaniu \texttt{std::function} w każdym przypadku.

W C++20 przechowywany obiekt docelowy musi być kopiowalny, dlatego
następujący kod nie jest poprawny:

\begin{lstlisting}
#include <functional>
#include <memory>

// ERROR in C++20: the closure is not copy-constructible.
// std::function<int()> f =
//     [p = std::make_unique<int>(9)] { return *p; };
\end{lstlisting}

Od \textbf{C++23} biblioteka udostępnia
\texttt{std::move\_only\_function}, która służy do przechowywania
nienadających się do kopiowania obiektów wywoływalnych. W C++20
można natomiast przechowywać taki obiekt bezpośrednio jako \texttt{auto}
lub przekazać go do szablonu, który nie wymaga kopiowania.

\section{Co wybrać: \texttt{auto}, szablon, wskaźnik czy \texttt{std::function}?}
\label{sec:lambda-choice}

\begin{center}
\begin{tabular}{|p{3.0cm}|p{4.4cm}|p{5.0cm}|}
\hline
\textbf{Forma} & \textbf{Zalety} & \textbf{Ograniczenia} \\
\hline
\texttt{auto f = lambda} & Konkretny typ, bez konieczności wymazywania typu, naturalna obsługa stanu & Zmienna ma ustalony typ domknięcia \\
\hline
\texttt{template<class F>} & Działa z wieloma typami, umożliwia inline i forwarding & Definicja szablonu zwykle dostępna w nagłówku; możliwe powiększanie kodu \\
\hline
Wskaźnik do funkcji & Prosty nieposiadający stanu interfejs & Nie przyjmuje lambd z przechwyceniami \\
\hline
\texttt{std::function} & Jedna sygnatura, zmienne zachowanie w czasie wykonania & Type erasure; koszt pośredniego wywołania; wymóg kopiowalności celu w C++20 \\
\hline
\end{tabular}
\end{center}

\paragraph{Praktyczna reguła.}
Jeżeli callback wykorzystujemy natychmiast, wybierajmy zwykle szablon
lub \texttt{auto}. Jeżeli musimy przechowywać i później podmieniać różne
callbacki o jednej sygnaturze, \texttt{std::function} jest często
czytelniejsza niż ręczne modelowanie polimorfizmu.

\section{Lambdy ze standardowymi algorytmami}
\label{sec:lambda-algorithms}

\subsection{\texttt{find\_if}, \texttt{count\_if}, \texttt{transform}}

\begin{lstlisting}
#include <algorithm>
#include <vector>

std::vector<int> xs{1, 2, 3, 4, 5, 6};

int limit = 3;
auto it = std::find_if(xs.begin(), xs.end(),
                       [limit](int x) { return x > limit; });
// *it == 4

auto even_count = std::count_if(xs.begin(), xs.end(),
                                [](int x) { return x % 2 == 0; });
// even_count == 3

std::vector<int> squares(xs.size());
std::transform(xs.begin(), xs.end(), squares.begin(),
               [](int x) { return x * x; });
\end{lstlisting}

Dla \texttt{std::find\_if} trzeba przed dereferencją sprawdzić, czy
wynik jest różny od \texttt{xs.end()}. Dla \texttt{std::transform} zbiór
docelowy musi już zawierać wystarczającą liczbę elementów: samo
\texttt{reserve()} zwiększa \texttt{capacity()}, lecz nie \texttt{size()}.
Alternatywą jest pusty wektor i iterator \texttt{std::back\_inserter}.

\subsection{\texttt{std::ranges} i projekcje (C++20)}

\begin{lstlisting}
#include <algorithm>
#include <ranges>
#include <string>
#include <vector>

struct Person {
    std::string name;
    int age;
};

std::vector<Person> people{{"Ada", 31}, {"Bob", 18}, {"Eve", 25}};

std::ranges::sort(people, {}, &Person::age);
// sorted by age ascending

auto adults = people | std::views::filter(
    [](const Person& p) { return p.age >= 18; });
\end{lstlisting}

Trzeci argument \texttt{ranges::sort} to \emph{projekcja}:
\texttt{\&Person::age} pobiera klucz sortowania, a domyślny komparator
porównuje klucze. \texttt{views::filter} nie buduje automatycznie nowego
wektora; jest leniwym widokiem, który odwołuje się do danych źródłowych.
Przy widokach szczególnie uważnie trzeba kontrolować czas życia
kontenera i przechwyceń predykatu.

\section{\texttt{std::ref}, \texttt{std::bind\_front} i częściowe wiązanie argumentów}
\label{sec:bind-front}

\begin{lstlisting}
#include <functional>

int add(int x, int y) { return x + y; }

auto add_ten = std::bind_front(add, 10); // C++20
int x = add_ten(5);                      // 15
\end{lstlisting}

\texttt{std::bind\_front} przechowuje funkcję i ustalone argumenty,
a pozostałe przyjmuje podczas wywołania. W wielu zastosowaniach
równoważna lambda jest czytelniejsza:

\begin{lstlisting}
auto add_ten_lambda = [](int x) { return add(10, x); };
\end{lstlisting}

Domyślnie argumenty przechowywane przez \texttt{bind\_front} są
przechowywane przez wartość, z uwzględnieniem przenoszenia. Jeśli
potrzebujemy odwołania do oryginalnego obiektu, stosujemy
\texttt{std::ref} (a dla dostępu tylko do odczytu
\texttt{std::cref}), pamiętając o czasie jego życia:

\begin{lstlisting}
#include <functional>

struct Accumulator {
    int total = 0;
    void add(int n) { total += n; }
};

Accumulator acc;
auto update = std::bind_front(&Accumulator::add, std::ref(acc));
update(7); // acc.total == 7
\end{lstlisting}

Starsze \texttt{std::bind} jest dostępne, ale ma bardziej skomplikowane
reguły dotyczące argumentów i podstawiania. W nowym kodzie zwykle
preferujemy lambdy lub \texttt{std::bind\_front}.

\section{Praktyczny projekt: rejestr callbacków}
\label{sec:callbacks}

Rozważmy klasę powiadamiającą o zmianie wartości. Do przechowywania
procedury obsługi zdarzenia użyjemy \texttt{std::function}. Dzięki
temu można ją zastąpić inną nawet po utworzeniu obiektu.

\begin{lstlisting}
#include <functional>
#include <utility>

class Counter {
public:
    using Callback = std::function<void(int)>;

    void on_change(Callback callback) {
        callback_ = std::move(callback);
    }

    void increment() {
        ++value_;
        if (callback_) callback_(value_);
    }

    int value() const noexcept { return value_; }

private:
    int value_ = 0;
    Callback callback_;
};

Counter counter;
counter.on_change([](int value) {
    // report the new value
});
counter.increment();
\end{lstlisting}

Kod w pojedynczym wątku jest prosty i wystarczający dla wielu
interfejsów zdarzeniowych. W programie wielowątkowym równoczesna
zmiana \texttt{callback\_} oraz wykonywanie \texttt{increment()}
wymaga dodatkowej synchronizacji. \texttt{std::function} nie zapewnia
automatycznie ochrony wewnętrznego stanu przed wyścigami danych.

\subsection{Callback a czas życia właściciela}

Typowy przebieg to: rejestrujemy callback, właściciel danych zostaje
zniszczony, a ktoś później wywołuje callback. Jeśli lambda zachowała
samą referencję lub surowy wskaźnik, dostęp byłby wtedy błędny.
Słaby wskaźnik w poniższym przykładzie pozwala sprawdzić czas życia.

\begin{lstlisting}
#include <memory>

struct Worker : std::enable_shared_from_this<Worker> {
    int total = 0;

    auto callback() {
        std::weak_ptr<Worker> weak = shared_from_this();
        return [weak](int n) {
            if (auto self = weak.lock()) {
                self->total += n;
            }
        };
    }
};
\end{lstlisting}

Przykład zakłada, że obiekt \texttt{Worker} jest już zarządzany przez
\texttt{std::shared\_ptr} w chwili wywołania \texttt{callback()}.
\texttt{weak\_ptr} pozwala sprawdzić, czy obiekt wciąż istnieje,
bez przedłużania jego życia przez przechwycenie silnego wskaźnika.
To ważne również dla unikania cykli własności, np. gdy sam obiekt
przechowuje callback odnoszący się do niego.

\section{Rekurencyjne lambdy bez \texttt{std::function}}

Lambda zwykle nie zna bezpośrednio własnej nazwy podczas definiowania.
Można jednak przekazywać obiekt wywoływalny jako argument:

\begin{lstlisting}
auto factorial = [](auto&& self, unsigned n) -> unsigned {
    return n <= 1 ? 1 : n * self(self, n - 1);
};

unsigned x = factorial(factorial, 5); // 120
\end{lstlisting}

To odmiana wzorca \emph{self-parameter}. Dla realnych zastosowań
rekurencyjnych warto ponadto zwrócić uwagę na przepełnienie liczb
całkowitych i maksymalną głębokość stosu. W C++20 metoda ta pozwala
uniknąć pośredniej alokacji i wymazywania typu, których mogłaby wymagać
rekurencja oparta na \texttt{std::function}.

\section{Lambdy w metaprogramowaniu: \texttt{std::visit} i \texttt{std::apply}}
\label{sec:lambda-visit-apply}

Dwa przydatne zastosowania łączą lambdy z mechanizmami z
rozdziałów~\ref{chap:cpp14-type-erasure} i~\ref{chap:variadic-templates-cpp}.
\texttt{std::visit} wybiera wywołanie na podstawie aktywnego typu
\texttt{std::variant}, natomiast \texttt{std::apply} rozwija argumenty
z krotki lub podobnej struktury.

\subsection{Zestaw przeciążeń oparty na dziedziczeniu}

\begin{lstlisting}
#include <string>
#include <variant>

template <class... Fs>
struct overloaded : Fs... {
    using Fs::operator()...;
};

std::variant<int, std::string> value = std::string{"Hello"};

auto stringify = overloaded{
    [](int x) { return std::to_string(x); },
    [](const std::string& x) { return x; }
};

std::string result = std::visit(stringify, value); // "Hello"
\end{lstlisting}

Każda lambda dostarcza jeden \texttt{operator()}. Dziedziczenie
wariadyczne i deklaracja \texttt{using Fs::operator()...} tworzą zbiór
przeciążeń, spośród których kompilator wybiera odpowiednią funkcję.
W tym przykładzie wszystkie warianty wywołania zwracają ten sam typ,
\texttt{std::string}, czego wymaga tu pojedyncze wywołanie
\texttt{std::visit}. Konstrukcja jest przydatna przy przetwarzaniu
stanów programu i heterogenicznych danych.

\subsection{Rozwijanie krotki przez \texttt{std::apply}}

\begin{lstlisting}
#include <tuple>

std::tuple<int, int, int> data{2, 3, 5};
int sum = std::apply(
    [](auto... values) { return (0 + ... + values); },
    data);
// sum == 10
\end{lstlisting}

Tutaj \texttt{std::apply} bierze elementy krotki i wywołuje lambdę
jak zwykłą funkcję trójargumentową. Sama lambda wykorzystuje
\emph{fold expression}, a więc nie trzeba pisać rekurencyjnego kodu.

\section{Pakiety przechwyceń i obliczenia odroczone (C++20)}

C++20 pozwala stosować \emph{pack expansion} również w przechwytywaniu
inicjalizującym. Możemy w ten sposób zbudować jeden obiekt domknięcia,
który przechowuje cały pakiet wartości do późniejszego przetworzenia:

\begin{lstlisting}
#include <concepts>

template <std::integral... Ts>
auto delayed_sum(Ts... xs) {
    return [...stored = xs] {
        return (0 + ... + stored);
    };
}

auto job = delayed_sum(2, 3, 5);
int result = job(); // 10
\end{lstlisting}

Lambdę można zwrócić z funkcji, ponieważ wartości są przechowywane
wewnątrz domknięcia. Przechwycenie przez referencję miałoby inne
wymagania dotyczące czasu życia. W rzeczywistych fabrykach callbacków
pakiet może zawierać obiekty niekopiowalne; wtedy należy zadbać
o przekazanie własności (np. przez \texttt{std::move} i
\texttt{std::forward}) oraz o wymagania interfejsu, który callback
przechowuje.

\section{Lambdy w obliczeniach stałych}

Jeśli ciało lambdy dopuszcza obliczanie stałe, jej wywołanie może
być częścią wyrażenia stałego. Od C++17 odpowiedni operator wywołania
jest domyślnie \texttt{constexpr}, gdy spełnia wymagania takiej funkcji:

\begin{lstlisting}
constexpr auto cube = [](int x) { return x * x * x; };
static_assert(cube(3) == 27);

// C++20: each call must be evaluated at compile time.
auto fourth_power = [](int x) consteval {
    return x * x * x * x;
};
static_assert(fourth_power(2) == 16);
\end{lstlisting}

\texttt{consteval} jest mocniejszą deklaracją niż \texttt{constexpr}:
wymaga, aby właściwe wywołania funkcji natychmiastowej spełniały
reguły obliczania stałego. Nie należy traktować tego jako gwarancji,
że każda zwykła lambda zostanie wykonana podczas kompilacji.

\section{Koncepcje dla obiektów wywoływalnych}
\label{sec:callable-concepts}

W C++20 można ograniczać typy wywoływalne m.in. następującymi
koncepcjami:
\begin{itemize}
\item \texttt{std::invocable} -- poprawność wywołania z danymi argumentami;
\item \texttt{std::regular\_invocable} -- dodatkowe wymagania semantyczne;
\item \texttt{std::predicate} -- poprawnie wywoływalny predykat z
wynikiem dającym się traktować jako wartość logiczną.
\end{itemize}
Koncepcje nie dowodzą wszystkich własności semantycznych. Przykładowo
kompilator nie zweryfikuje matematycznie stabilności wyników predykatu.

\begin{lstlisting}
#include <concepts>
#include <functional>
#include <utility>

template <class F, class T>
    requires std::predicate<F&, const T&>
bool test(const T& value, F&& predicate) {
    return std::invoke(std::forward<F>(predicate), value);
}

bool ok = test(8, [](int x) { return x % 2 == 0; });
\end{lstlisting}

\textbf{Istotne rozróżnienie:} koncepcja weryfikuje przede wszystkim
własności typów i poprawność wyrażeń. Nie dowodzi matematycznie
własności predykatu, np. że comparator tworzy ścisły słaby porządek.
Programista nadal odpowiada za kontrakt semantyczny.

\section{Typowe błędy i trudniejsze przypadki}
\label{sec:lambda-pitfalls}

\subsection{Referencja przechwycona do operacji asynchronicznej}

\begin{lstlisting}
#include <future>

std::future<int> start() {
    int x = 7;
    // WRONG: the task could outlive x.
    // return std::async(std::launch::async, [&x] { return x; });
    return std::async(std::launch::async, [x] { return x; });
}
\end{lstlisting}

W zadaniach asynchronicznych preferuj przechwycenia przez wartość
lub świadomie zarządzaną własność, dopóki nie mamy ścisłej gwarancji,
że obiekt referowany przetrwa wykonanie zadania. Sama biblioteka
\texttt{std::async} nie wydłuża życia lokalnych zmiennych funkcji
uruchamiającej.

\subsection{Predykat sortowania nie może naruszać porządku}

\begin{lstlisting}
// WRONG: <= is not a strict ordering relation.
// std::sort(xs.begin(), xs.end(),
//           [](int a, int b) { return a <= b; });

std::sort(xs.begin(), xs.end(),
          [](int a, int b) { return a < b; });
\end{lstlisting}

Komparator do \texttt{sort} musi definiować \emph{ścisły słaby
porządek}. Przykładowo \texttt{comp(x,x)} ma dawać \texttt{false}.
Naruszenie wymagań algorytmu może prowadzić do niezdefiniowanego
zachowania, mimo że typy argumentów się zgadzają.

\subsection{Pomylenie kopii z oryginalnym stanem}

\begin{lstlisting}
int sum = 0;
std::for_each(xs.begin(), xs.end(), [sum](int x) mutable {
    sum += x; // modifies only the closure's copy
});
// outer sum remains 0

std::for_each(xs.begin(), xs.end(), [&sum](int x) {
    sum += x;
});
\end{lstlisting}

Nie należy zakładać, że \texttt{mutable} zmienia zewnętrzną
zmienną. Należy też pamiętać, że algorytmy
mogą kopiować obiekt funkcyjny; nie opierajmy semantyki algorytmu
na nieudokumentowanych szczegółach kopiowania predykatu.

\subsection{Przechwycenie wskaźnika nie przedłuża życia obiektu}

\begin{lstlisting}
int* p = /* pointer to some object */ nullptr;
auto f = [p] {
    if (p) return *p;
    return 0;
};
\end{lstlisting}

Wartość wskaźnika może zostać skopiowana do lambdy, ale wskazywany
obiekt nadal musi istnieć w chwili dereferencji. Przechwycenie
\texttt{shared\_ptr} przenosi własność współdzieloną, a
\texttt{weak\_ptr} pozwala ją warunkowo odzyskać. To nie jest to
samo co skopiowanie wskaźnika surowego.

\section{Zestawienie: kiedy stosować którą konstrukcję?}

\begin{center}
\begin{tabular}{|p{5.7cm}|p{6.6cm}|}
\hline
\textbf{Potrzeba} & \textbf{Zalecany wybór} \\
\hline
Krótki predykat do algorytmu & Lambda z jawnymi przechwyceniami \\
\hline
Wywołanie bez stanu w API C & Zwykła funkcja albo zgodna lambda bez przechwyceń jako wskaźnik do funkcji \\
\hline
Wydajny callback o typie znanym w czasie kompilacji & Parametr szablonowy \texttt{F}, ewentualnie \texttt{std::invocable} \\
\hline
Callback z możliwością późniejszej podmiany & \texttt{std::function} (C++20, kopiowalny cel) \\
\hline
Przechowanie \texttt{unique\_ptr} wewnątrz lambdy & Init-capture i przeniesienie; nie przechowuj takiej lambdy w \texttt{std::function} w C++20 \\
\hline
Callback żyjący dłużej niż zmienne lokalne & Kopiuj potrzebny stan lub zarządzaj życiem przez RAII i smart pointery \\
\hline
Uniwersalne wywołanie dowolnego callable & \texttt{std::invoke} i \texttt{std::forward} \\
\hline
Częściowe podstawienie kilku argumentów & Lambda lub \texttt{std::bind\_front} \\
\hline
Bezpieczny predykat dla szablonu & \texttt{std::predicate} plus kontrola jego semantyki \\
\hline
\end{tabular}
\end{center}

\paragraph{Pytania kontrolne.}
\begin{enumerate}
\item Dlaczego dwie lambdy o identycznym kodzie mają różne typy?
\item Jaka jest różnica między \texttt{[x]}, \texttt{[\&x]} i \texttt{[x]() mutable}?
\item Kiedy bezpiecznie zwrócić lambdę przechwytującą zmienną lokalną?
\item Dlaczego lambda posiadająca \texttt{unique\_ptr} nie może być celem \texttt{std::function} w C++20?
\item Jak \texttt{std::invoke} wywołuje wskaźnik do metody klasy?
\item Czym różnią się \texttt{[this]} i \texttt{[*this]}?
\item Co może pójść źle, gdy \texttt{std::views::filter} korzysta z lokalnej zmiennej przez referencję?
\item Kiedy interfejs przyjmujący szablon \texttt{F} jest lepszy od \texttt{std::function}?
\item Dlaczego komparator \texttt{a <= b} nie jest poprawny dla \texttt{std::sort}?
\item W jakiej sytuacji przechwycenie \texttt{weak\_ptr} jest lepsze od \texttt{shared\_ptr}?
\end{enumerate}

\paragraph{Dalsza lektura.}
W dokumentacji \emph{cppreference} warto przeczytać hasła
\emph{Lambda expressions}, \emph{Function objects} i \emph{Ranges}.
Dokładne opisy narzędzi znajdują się również pod hasłami:
\begin{itemize}
\item \texttt{std::invoke} oraz \texttt{std::function};
\item \texttt{std::bind\_front} oraz \texttt{std::predicate}.
\end{itemize}


\chapter[Lambdy i obiekty funkcyjne]{Wyrażenia lambda, obiekty funkcyjne i abstrakcje wywoływalne}
\label{chap:lambdas-cpp}

\section{Po co istnieją obiekty funkcyjne?}

W C++ funkcję można reprezentować nie tylko przez jej nazwę lub wskaźnik.
Można również utworzyć \emph{obiekt wywoływalny} (\textit{callable object}):
obiekt, który daje się zastosować w wyrażeniu przypominającym wywołanie
funkcji. Pozwala to przekazywać zachowanie jako parametr algorytmu,
zapamiętywać stan pomiędzy wywołaniami oraz tworzyć elastyczne API
(np. sortowanie z własnym porządkiem, filtrowanie, obsługa zdarzeń).

\begin{lstlisting}
#include <algorithm>
#include <vector>

std::vector<int> numbers{5, 2, 8, 1};
std::sort(numbers.begin(), numbers.end(),
          [](int a, int b) { return a > b; });
// numbers == {8, 5, 2, 1}
\end{lstlisting}

Trzeci argument \texttt{std::sort} to reguła porównywania. Zamiast pisać
osobną funkcję o globalnej nazwie, definiujemy ją dokładnie w miejscu
użycia. Wyrażenie \texttt{[](int a, int b) \{ return a > b; \}}
jest \emph{lambdą}: tworzy obiekt unikatowego typu domknięcia
(\textit{closure type}).

\paragraph{Zakres i wymagania.}
Przyjmujemy \textbf{C++20}. Lambdy wprowadzono w C++11, przechwytywanie
inicjalizujące i generyczne lambdy w C++14, przechwytywanie kopii
\texttt{*this} w C++17, a jawne listy parametrów szablonu lambdy
w C++20. Wykorzystamy semantykę przenoszenia z rozdziału~\ref{chap:move-semantics-cpp},
szablony wariadyczne z rozdziału~\ref{chap:variadic-templates-cpp}
i koncepcje z rozdziału~\ref{chap:concepts-cpp}.

\section{Trzy podstawowe postacie kodu wywoływalnego}
\label{sec:callables-forms}

\subsection{Funkcja i wskaźnik do funkcji}

Zwykła funkcja nie przechowuje kontekstu wywołania. Jej adres można
przypisać wskaźnikowi o odpowiednim typie:

\begin{lstlisting}
int twice(int x) { return 2 * x; }

int (*fp)(int) = &twice;
int a = fp(7);       // 14
int b = (*fp)(7);    // also 14
\end{lstlisting}

W deklaracji \texttt{int (*fp)(int)} nawiasy są istotne. Sam wskaźnik
przechowuje adres funkcji; nie ma miejsca na dodatkowy stan, taki jak
współczynnik mnożenia ustalany w czasie działania programu.

\subsection{Funktor: klasa z \texttt{operator()}}

\begin{lstlisting}
struct Multiply {
    int factor;

    int operator()(int x) const {
        return factor * x;
    }
};

Multiply times_three{3};
int result = times_three(7); // 21
\end{lstlisting}

\texttt{Multiply} jest klasą, a \texttt{times\_three} jej egzemplarzem.
Dzięki polu \texttt{factor} każdy egzemplarz może reprezentować inne
zachowanie. Operator wywołania może być przeciążany, szablonowy,
opatrzony kwalifikatorami \texttt{const} lub \texttt{\&\&} i mieć
złożoną implementację.

\subsection{Lambda jako wygodny zapis funktora}

\begin{lstlisting}
int factor = 3;
auto times_three = [factor](int x) {
    return factor * x;
};
int result = times_three(7); // 21
\end{lstlisting}

Kompilator tworzy odrębny, nienazwany typ klasy odpowiadający lambdzie.
Uproszczony model jej działania jest następujący:

\begin{lstlisting}
struct ClosureModel {
    int factor;

    int operator()(int x) const {
        return factor * x;
    }
};
\end{lstlisting}

Jest to \emph{model pojęciowy}, nie obiecana przez standard reprezentacja
pamięci. Każde wystąpienie wyrażenia lambda ma własny typ, nawet jeżeli
zapis kodu dwóch lambd wygląda identycznie.

\section{Budowa wyrażenia lambda}
\label{sec:lambda-syntax}

Najważniejsze elementy zapisu przedstawia schemat:

\begin{lstlisting}
[captures](parameters) mutable noexcept -> ResultType {
    // body
}
\end{lstlisting}

W C++20 można również jawnie określić parametry szablonu, koncepcje
oraz klauzulę \texttt{requires}. Typ wyniku zwykle jest wywnioskowany
z instrukcji \texttt{return}; podanie \texttt{-> ResultType} bywa
przydatne, gdy chcemy wymusić konkretny typ lub zwracać referencję.

\begin{lstlisting}
auto square = [](int x) { return x * x; };
auto half = [](double x) -> double { return x / 2; };
auto truth = [] { return true; }; // () may be omitted

int a = square(4);  // 16
double b = half(5); // 2.5
\end{lstlisting}

\paragraph{Najważniejsze zastrzeżenie.}
Lambda nie jest wskaźnikiem do funkcji. Zwykle tworzy obiekt z
\texttt{operator()}; jedynie pewne lambdy bez przechwyceń mają
konwersję na wskaźnik do funkcji o pasującej sygnaturze.

\section{Przechwytywanie zmiennych (\textit{captures})}

\begin{center}\small
\begin{tabularx}{\textwidth}{@{}p{0.30\textwidth}X@{}}
\toprule
Zapis lambdy & Znaczenie \\
\midrule
\code{[x]} & Przechwytuje wartość \code{x} przez kopię w obiekcie domknięcia. \\
\texttt{[\&x]} & Przechwytuje referencję do \code{x}; wywołanie wymaga, by \code{x} nadal żyło. \\
\code{[=]} & Domyślnie przechwytuje użyte zmienne lokalne przez wartość. \\
\texttt{[\&]} & Domyślnie przechwytuje użyte zmienne lokalne przez referencję. \\
\code{[p = std::move(p)]} & Przenosi wartość do składowej domknięcia przez inicjalizację przechwycenia. \\
\code{mutable} & Pozwala zmieniać kopie przechwycone przez wartość wewnątrz lambdy. \\
\bottomrule
\end{tabularx}
\end{center}
\label{sec:lambda-captures}

\subsection{Przechwycenie przez wartość}

\begin{lstlisting}
int threshold = 10;
auto above = [threshold](int x) {
    return x > threshold;
};
auto above_current = [&threshold](int x) {
    return x > threshold;
};
threshold = 100;

bool first = above(20);         // true: zapamietane 10
bool second = above_current(20); // false: aktualne 100
\end{lstlisting}

\texttt{[threshold]} inicjalizuje składową obiektu domknięcia wartością
\texttt{threshold} w chwili tworzenia lambdy. Późniejsza zmiana zmiennej
zewnętrznej nie modyfikuje tej kopii. Jeśli kopiowany obiekt ma kosztowne
kopiowanie, samo utworzenie lub skopiowanie lambdy może także mieć koszt.

\subsection{Przechwycenie przez referencję}

\begin{lstlisting}
int total = 0;
auto add = [&total](int value) {
    total += value;
};
add(4);
add(8); // total == 12
\end{lstlisting}

\texttt{[\&total]} uzyskuje dostęp do tej samej zmiennej, nie do jej
kopii. Lambda \textbf{nie wydłuża czasu życia} obiektu przechwyconego
przez referencję. Dlatego jest wygodna, gdy callback wykonywany jest
synchronicznie, ale potencjalnie niebezpieczna, jeśli jest zapamiętywany
albo wykonywany później.

\subsection{Przechwycenia domyślne i mieszane}

\begin{lstlisting}
int a = 2;
int b = 3;
int count = 0;

auto first  = [=]  { return a + b; }; // by copy
auto second = [&]  { ++count; };     // by reference
auto third  = [=, &count] {
    ++count;
    return a + b;
};
\end{lstlisting}

\texttt{[=]} i \texttt{[\&]} obejmują te użyte w treści zmienne lokalne,
które podlegają niejawnemu przechwyceniu; nie oznaczają automatycznego
skopiowania \emph{wszystkich} zmiennych z otoczenia. W większym kodzie
zwykle lepiej wypisać przechwycenia jawnie: dzięki temu widać zależności
i łatwiej ocenić bezpieczeństwo czasu życia.

\subsection{Przechwytywanie inicjalizujące (C++14)}

W nawiasach kwadratowych można zdefiniować nową nazwę i sposób jej
inicjalizacji. Dzięki temu lambdy mogą przechowywać wartości pochodne,
a także zasoby, których nie wolno kopiować:

\begin{lstlisting}
#include <memory>
#include <utility>

auto p = std::make_unique<int>(42);
auto read = [q = std::move(p)] {
    return *q;
};
// p == nullptr; ownership is now in the closure
int answer = read(); // 42
\end{lstlisting}

Taka lambda jest \emph{move-only}, ponieważ zawiera
\texttt{std::unique\_ptr}. Nie można jej skopiować. To istotne przy
wyborze miejsca przechowywania callbacków: w C++20
\texttt{std::function} wymaga kopiowalnego obiektu docelowego.

\subsection{Niebezpieczne przechwytywanie: wisząca referencja}

\begin{lstlisting}
// WRONG: do not call the returned lambda!
auto broken() {
    int local = 5;
    return [&local] { return local; };
}

// Correct: capture the value.
auto safe() {
    int local = 5;
    return [local] { return local; };
}
\end{lstlisting}

\texttt{broken()} zwraca obiekt z dostępem do zmiennej, która przestaje
istnieć wraz z wyjściem z funkcji. Późniejsze wywołanie prowadziłoby
do niezdefiniowanego zachowania (\textit{undefined behavior}).
Przechwycenie przez kopię rozwiązuje ten konkretny problem.

\section{Przechwytywanie \texttt{this} i \texttt{*this}}
\label{sec:lambda-this}

W metodzie klasy można korzystać z jej pól wewnątrz lambdy. Trzeba
jednak rozróżnić przechwycenie wskaźnika \texttt{this} od skopiowania
bieżącego obiektu.

\begin{lstlisting}
#include <string>

struct Label {
    std::string text;

    auto by_pointer() const {
        return [this] { return text; };
    }

    auto by_copy() const {
        return [*this] { return text; }; // C++17
    }
};
\end{lstlisting}

Lambda zwrócona przez \texttt{by\_pointer()} odwołuje się do oryginalnego
obiektu. Nie może być bezpiecznie wywoływana po jego zniszczeniu.
Lambda z \texttt{[*this]} przechowuje kopię obiektu (jeśli da się go
skopiować); dla składowych zarządzanych przez RAII jest to zwykle
bardziej przewidywalne, ale może być kosztowne.

\textbf{Uwaga:} zapis \texttt{[=]} użyty w metodzie nie kopiuje automatycznie
całego obiektu. Niejawne przechwytywanie \texttt{this} przez \texttt{[=]}
jest w C++20 przestarzałe (\textit{deprecated}); zapisujmy
\texttt{[this]} jawnie albo \texttt{[*this]}, gdy potrzebujemy kopii.

\section{Domyślnie stały operator wywołania i \texttt{mutable}}

Dla typowej lambdy bez słowa \texttt{mutable} operator wywołania ma
kwalifikator \texttt{const}. Oznacza to, że przechwyconych przez wartość
składowych nie wolno w nim zwyczajnie modyfikować:

\begin{lstlisting}
int counter = 0;

// ERROR: captured counter is not writable in this closure.
// auto bad = [counter] { ++counter; };

auto next = [counter]() mutable {
    return ++counter;
};

int a = next(); // 1
int b = next(); // 2
// outer counter is still 0
\end{lstlisting}

\texttt{mutable} modyfikuje zachowanie \emph{kopii wewnątrz lambdy},
nie powoduje przechwycenia oryginalnej zmiennej przez referencję.
Różne kopie obiektu domknięcia mogą przechowywać różny stan:

\begin{lstlisting}
auto a = [n = 0]() mutable { return ++n; };
auto b = a; // independent copy of state

a(); // 1
a(); // 2
b(); // 1
\end{lstlisting}

Przy przechwyceniu przez referencję do modyfikacji zewnętrznego
obiektu \texttt{mutable} nie jest potrzebne: stałość metody
\texttt{operator()} nie zmienia stałości obiektu osiąganego przez
referencję.

\section{Typ lambdy, kopiowanie i konwersja na wskaźnik}

\begin{lstlisting}
#include <type_traits>

inline constexpr auto f = [](int x) { return x + 1; };

auto a = [](int x) { return x; };
auto b = [](int x) { return x; };
static_assert(!std::is_same_v<decltype(a), decltype(b)>);

int (*pf)(int) = +[](int x) { return x + 1; };
\end{lstlisting}

Dwie osobno zapisane lambdy mają różne typy, nawet jeśli robią to samo.
W C++20 typ domknięcia lambdy bez przechwyceń jest domyślnie
konstruowalny; dostępność różnych operacji kopiowania/przypisywania
zależy od tego, czy lambda przechwytuje zmienne. Dla lambd z
przechwyceniami obowiązują też ograniczenia narzucone przez typy pól
(na przykład niekopiowalność \texttt{unique\_ptr}).

\texttt{+[]} jest znanym idiomem wymuszającym konwersję
nieprzechwytującej żadnej zmiennej lambdy na wskaźnik do funkcji.
Nie działa dla lambd przechwytujących:

\begin{lstlisting}
int offset = 10;
auto with_state = [offset](int x) { return x + offset; };
// int (*bad)(int) = with_state;  // ERROR
\end{lstlisting}

\section{Lambdy generyczne i własne parametry szablonu}
\label{sec:lambda-generic}

\subsection{Parametry \texttt{auto} (C++14)}

\begin{lstlisting}
#include <string>
auto print_twice = [](const auto& value) {
    return value + value;
};

auto x = print_twice(7);                 // 14
auto s = print_twice(std::string("ab")); // "abab"
\end{lstlisting}

Parametr \texttt{auto} sprawia, że \texttt{operator()} obiektu domknięcia
jest szablonem funkcji. Przy wywołaniu kompilator dedukuje odpowiednie
typy tak jak przy zwykłej funkcji szablonowej.

\subsection{Jawne parametry szablonu w lambdzie (C++20)}

\begin{lstlisting}
#include <concepts>

auto equal_types = []<class T>(const T& a, const T& b) {
    return a == b;
};

auto sum_integrals = []<std::integral... Ts>(Ts... values) {
    return (0 + ... + values);
};

bool yes = equal_types(3, 3);
int total = sum_integrals(1, 2, 3, 4); // 10
\end{lstlisting}

W przypadku \texttt{equal\_types} oba argumenty muszą dać ten sam
dedukowany typ \texttt{T}. W drugim przykładzie łączymy \emph{concepts}
z pakietem parametrów i wyrażeniem fold znanym z
rozdziałów~\ref{chap:variadic-templates-cpp} i~\ref{chap:concepts-cpp}.

\subsection{Uniwersalny wrapper z perfect forwarding}

\begin{lstlisting}
#include <functional>
#include <utility>

auto invoke_once = []<class F, class... Args>(
    F&& f, Args&&... args) -> decltype(auto) {
    return std::invoke(std::forward<F>(f),
                       std::forward<Args>(args)...);
};
\end{lstlisting}

\texttt{F\&\&} i \texttt{Args\&\&...} są tu referencjami przekazującymi,
a \texttt{decltype(auto)} zachowuje typ wyniku, w tym jego
referencyjność. W odróżnieniu od \texttt{std::move},
\texttt{std::forward} nie wymusza kategorii rvalue dla każdego argumentu.
Dzięki \texttt{std::invoke} wrapper działa nie tylko dla zwykłych
funkcji i lambd, lecz także dla wskaźników do składowych.

\section{Wymuszanie sygnatury: \texttt{noexcept}, \texttt{requires}, zwracanie referencji}

\begin{lstlisting}
#include <concepts>
#include <utility>

auto positive = [](int x) noexcept {
    return x > 0;
};

auto increment = []<std::integral T>(T x) -> T {
    return x + 1;
};

auto first = []<class T>(T& x) -> T& {
    return x;
};

int value = 3;
first(value) = 17; // modifies value
\end{lstlisting}

Ważne są dwie subtelności. Po pierwsze, \texttt{noexcept} jest obietnicą
programisty; jeśli wyjątek wydostanie się z funkcji oznaczonej
\texttt{noexcept}, wywołane zostanie \texttt{std::terminate}.
Po drugie, domyślna dedukcja wyniku lambdy deklarowanej bez
\texttt{decltype(auto)} lub jawnego \texttt{-> T\&} zwykle usuwa
referencję. Właściwe zwracanie referencji wymaga zatem świadomego
określenia typu wyniku i zapewnienia jego poprawnego czasu życia.

\section{\texttt{std::invoke}: jeden mechanizm wywołania}

\begin{center}\small
\begin{tabularx}{\textwidth}{@{}p{0.34\textwidth}X@{}}
\toprule
Narzędzie & Za co odpowiada \\
\midrule
\code{std::invoke(f, args...)} & Jednolicie wywołuje funkcję, lambdę lub wskaźnik do składowej. \\
\code{std::function<R(Args...)>} & Przechowuje kopiowalny obiekt wywoływalny o ustalonej sygnaturze przez wymazywanie typu. \\
\code{std::ref(x)} & Opakowuje referencję, gdy API normalnie kopiuje argument. \\
\code{std::bind_front(f, args...)} & Wiąże pierwsze argumenty funkcji i zwraca nowy obiekt wywoływalny. \\
\code{noexcept} przy lambdzie & Określa, że operator wywołania nie powinien pozwolić wyjątkowi opuścić funkcji. \\
\bottomrule
\end{tabularx}
\end{center}
\label{sec:std-invoke}

\texttt{std::invoke} z nagłówka \texttt{<functional>} wywołuje funkcję,
funktor lub lambdę, ale obsługuje także wskaźnik do metody i pola.

\begin{lstlisting}
#include <functional>

struct Box {
    int value = 10;
    int add(int x) const { return value + x; }
};

Box b;
int a = std::invoke([](int x) { return x * 2; }, 4);
int c = std::invoke(&Box::add, b, 7);      // 17
int& r = std::invoke(&Box::value, b);
r = 25;                                   // b.value == 25
\end{lstlisting}

To podstawowa cegiełka generycznych algorytmów. W nagłówku
\texttt{<type\_traits>} zdefiniowano m.in. \texttt{std::invoke\_result\_t},
\texttt{std::is\_invocable\_v} i \texttt{std::is\_nothrow\_invocable\_v},
które pozwalają sprawdzić w czasie kompilacji poprawność wywołania.

\begin{lstlisting}
#include <type_traits>

static_assert(std::is_invocable_v<decltype(positive), int>);
static_assert(!std::is_invocable_v<decltype(positive),
                                   std::string>);
static_assert(std::is_nothrow_invocable_v<decltype(positive),
                                          int>);
\end{lstlisting}

\section{\texttt{std::function}: jeden typ dla wielu zachowań}
\label{sec:std-function}

Lambdy o różnych typach nie mogą bezpośrednio zastępować się w jednej
zmiennej zadeklarowanej przez \texttt{auto}. Jeżeli chcemy przechowywać
\emph{dowolny} kopiowalny obiekt wywoływalny o zadanej sygnaturze,
możemy użyć \texttt{std::function<R(Args...)>}.

\begin{lstlisting}
#include <functional>
#include <string>

std::function<int(int)> operation =
    [](int x) { return 2 * x; };

int factor = 5;
operation = [factor](int x) { return factor * x; };
int result = operation(3); // 15

operation = nullptr;
if (!operation) {
    // no callable stored
}
\end{lstlisting}

\texttt{std::function} stosuje \emph{type erasure}: ukrywa rzeczywisty
typ obiektu docelowego za wspólnym interfejsem wywołania. Może
przechowywać funkcje, funktory i lambdy z przechwyceniami. Nie jest
odpowiednikiem \texttt{auto}: deklaracja \texttt{auto} zachowuje
konkretny typ i pozwala zwykle kompilatorowi łatwiej zoptymalizować
wywołanie.

\subsection{Pusty \texttt{std::function} i wyjątek}

\begin{lstlisting}
#include <functional>

std::function<void()> callback;
if (callback) {
    callback();
}
// Calling callback() while empty throws std::bad_function_call.
\end{lstlisting}

Sprawdzenie \texttt{if (callback)} jest zatem standardową częścią
obsługi opcjonalnych callbacków. Nie należy przechowywać
\texttt{std::function} bez potrzeby, jeśli konkretny typ wywoływalny
jest znany w czasie kompilacji.

\subsection{Koszty i ograniczenie kopiowalności}

Zwykle \texttt{std::function} wymaga pośredniego wywołania. Przechowywanie
większego domknięcia może też wymagać dynamicznej alokacji (zależy to
od implementacji i rozmiaru obiektu). To koszt elastyczności API, nie
argument przeciw używaniu \texttt{std::function} w każdym przypadku.

W C++20 przechowywany obiekt docelowy musi być kopiowalny, dlatego
następujący kod nie jest poprawny:

\begin{lstlisting}
#include <functional>
#include <memory>

// ERROR in C++20: the closure is not copy-constructible.
// std::function<int()> f =
//     [p = std::make_unique<int>(9)] { return *p; };
\end{lstlisting}

Od \textbf{C++23} biblioteka udostępnia
\texttt{std::move\_only\_function}, która służy do przechowywania
nienadających się do kopiowania obiektów wywoływalnych. W C++20
można natomiast przechowywać taki obiekt bezpośrednio jako \texttt{auto}
lub przekazać go do szablonu, który nie wymaga kopiowania.

\section{Co wybrać: \texttt{auto}, szablon, wskaźnik czy \texttt{std::function}?}
\label{sec:lambda-choice}

\begin{center}
\begin{tabular}{|p{3.0cm}|p{4.4cm}|p{5.0cm}|}
\hline
\textbf{Forma} & \textbf{Zalety} & \textbf{Ograniczenia} \\
\hline
\texttt{auto f = lambda} & Konkretny typ, bez konieczności wymazywania typu, naturalna obsługa stanu & Zmienna ma ustalony typ domknięcia \\
\hline
\texttt{template<class F>} & Działa z wieloma typami, umożliwia inline i forwarding & Definicja szablonu zwykle dostępna w nagłówku; możliwe powiększanie kodu \\
\hline
Wskaźnik do funkcji & Prosty nieposiadający stanu interfejs & Nie przyjmuje lambd z przechwyceniami \\
\hline
\texttt{std::function} & Jedna sygnatura, zmienne zachowanie w czasie wykonania & Type erasure; koszt pośredniego wywołania; wymóg kopiowalności celu w C++20 \\
\hline
\end{tabular}
\end{center}

\paragraph{Praktyczna reguła.}
Jeżeli callback wykorzystujemy natychmiast, wybierajmy zwykle szablon
lub \texttt{auto}. Jeżeli musimy przechowywać i później podmieniać różne
callbacki o jednej sygnaturze, \texttt{std::function} jest często
czytelniejsza niż ręczne modelowanie polimorfizmu.

\section{Lambdy ze standardowymi algorytmami}
\label{sec:lambda-algorithms}

\subsection{\texttt{find\_if}, \texttt{count\_if}, \texttt{transform}}

\begin{lstlisting}
#include <algorithm>
#include <vector>

std::vector<int> xs{1, 2, 3, 4, 5, 6};

int limit = 3;
auto it = std::find_if(xs.begin(), xs.end(),
                       [limit](int x) { return x > limit; });
// *it == 4

auto even_count = std::count_if(xs.begin(), xs.end(),
                                [](int x) { return x % 2 == 0; });
// even_count == 3

std::vector<int> squares(xs.size());
std::transform(xs.begin(), xs.end(), squares.begin(),
               [](int x) { return x * x; });
\end{lstlisting}

Dla \texttt{std::find\_if} trzeba przed dereferencją sprawdzić, czy
wynik jest różny od \texttt{xs.end()}. Dla \texttt{std::transform} zbiór
docelowy musi już zawierać wystarczającą liczbę elementów: samo
\texttt{reserve()} zwiększa \texttt{capacity()}, lecz nie \texttt{size()}.
Alternatywą jest pusty wektor i iterator \texttt{std::back\_inserter}.

\subsection{\texttt{std::ranges} i projekcje (C++20)}

\begin{lstlisting}
#include <algorithm>
#include <ranges>
#include <string>
#include <vector>

struct Person {
    std::string name;
    int age;
};

std::vector<Person> people{{"Ada", 31}, {"Bob", 18}, {"Eve", 25}};

std::ranges::sort(people, {}, &Person::age);
// sorted by age ascending

auto adults = people | std::views::filter(
    [](const Person& p) { return p.age >= 18; });
\end{lstlisting}

Trzeci argument \texttt{ranges::sort} to \emph{projekcja}:
\texttt{\&Person::age} pobiera klucz sortowania, a domyślny komparator
porównuje klucze. \texttt{views::filter} nie buduje automatycznie nowego
wektora; jest leniwym widokiem, który odwołuje się do danych źródłowych.
Przy widokach szczególnie uważnie trzeba kontrolować czas życia
kontenera i przechwyceń predykatu.

\section{\texttt{std::ref}, \texttt{std::bind\_front} i częściowe wiązanie argumentów}
\label{sec:bind-front}

\begin{lstlisting}
#include <functional>

int add(int x, int y) { return x + y; }

auto add_ten = std::bind_front(add, 10); // C++20
int x = add_ten(5);                      // 15
\end{lstlisting}

\texttt{std::bind\_front} przechowuje funkcję i ustalone argumenty,
a pozostałe przyjmuje podczas wywołania. W wielu zastosowaniach
równoważna lambda jest czytelniejsza:

\begin{lstlisting}
auto add_ten_lambda = [](int x) { return add(10, x); };
\end{lstlisting}

Domyślnie argumenty przechowywane przez \texttt{bind\_front} są
przechowywane przez wartość, z uwzględnieniem przenoszenia. Jeśli
potrzebujemy odwołania do oryginalnego obiektu, stosujemy
\texttt{std::ref} (a dla dostępu tylko do odczytu
\texttt{std::cref}), pamiętając o czasie jego życia:

\begin{lstlisting}
#include <functional>

struct Accumulator {
    int total = 0;
    void add(int n) { total += n; }
};

Accumulator acc;
auto update = std::bind_front(&Accumulator::add, std::ref(acc));
update(7); // acc.total == 7
\end{lstlisting}

Starsze \texttt{std::bind} jest dostępne, ale ma bardziej skomplikowane
reguły dotyczące argumentów i podstawiania. W nowym kodzie zwykle
preferujemy lambdy lub \texttt{std::bind\_front}.

\section{Praktyczny projekt: rejestr callbacków}
\label{sec:callbacks}

Rozważmy klasę powiadamiającą o zmianie wartości. Do przechowywania
procedury obsługi zdarzenia użyjemy \texttt{std::function}. Dzięki
temu można ją zastąpić inną nawet po utworzeniu obiektu.

\begin{lstlisting}
#include <functional>
#include <utility>

class Counter {
public:
    using Callback = std::function<void(int)>;

    void on_change(Callback callback) {
        callback_ = std::move(callback);
    }

    void increment() {
        ++value_;
        if (callback_) callback_(value_);
    }

    int value() const noexcept { return value_; }

private:
    int value_ = 0;
    Callback callback_;
};

Counter counter;
counter.on_change([](int value) {
    // report the new value
});
counter.increment();
\end{lstlisting}

Kod w pojedynczym wątku jest prosty i wystarczający dla wielu
interfejsów zdarzeniowych. W programie wielowątkowym równoczesna
zmiana \texttt{callback\_} oraz wykonywanie \texttt{increment()}
wymaga dodatkowej synchronizacji. \texttt{std::function} nie zapewnia
automatycznie ochrony wewnętrznego stanu przed wyścigami danych.

\subsection{Callback a czas życia właściciela}

Typowy przebieg to: rejestrujemy callback, właściciel danych zostaje
zniszczony, a ktoś później wywołuje callback. Jeśli lambda zachowała
samą referencję lub surowy wskaźnik, dostęp byłby wtedy błędny.
Słaby wskaźnik w poniższym przykładzie pozwala sprawdzić czas życia.

\begin{lstlisting}
#include <memory>

struct Worker : std::enable_shared_from_this<Worker> {
    int total = 0;

    auto callback() {
        std::weak_ptr<Worker> weak = shared_from_this();
        return [weak](int n) {
            if (auto self = weak.lock()) {
                self->total += n;
            }
        };
    }
};
\end{lstlisting}

Przykład zakłada, że obiekt \texttt{Worker} jest już zarządzany przez
\texttt{std::shared\_ptr} w chwili wywołania \texttt{callback()}.
\texttt{weak\_ptr} pozwala sprawdzić, czy obiekt wciąż istnieje,
bez przedłużania jego życia przez przechwycenie silnego wskaźnika.
To ważne również dla unikania cykli własności, np. gdy sam obiekt
przechowuje callback odnoszący się do niego.

\section{Rekurencyjne lambdy bez \texttt{std::function}}

Lambda zwykle nie zna bezpośrednio własnej nazwy podczas definiowania.
Można jednak przekazywać obiekt wywoływalny jako argument:

\begin{lstlisting}
auto factorial = [](auto&& self, unsigned n) -> unsigned {
    return n <= 1 ? 1 : n * self(self, n - 1);
};

unsigned x = factorial(factorial, 5); // 120
\end{lstlisting}

To odmiana wzorca \emph{self-parameter}. Dla realnych zastosowań
rekurencyjnych warto ponadto zwrócić uwagę na przepełnienie liczb
całkowitych i maksymalną głębokość stosu. W C++20 metoda ta pozwala
uniknąć pośredniej alokacji i wymazywania typu, których mogłaby wymagać
rekurencja oparta na \texttt{std::function}.

\section{Lambdy w metaprogramowaniu: \texttt{std::visit} i \texttt{std::apply}}
\label{sec:lambda-visit-apply}

Dwa przydatne zastosowania łączą lambdy z mechanizmami z
rozdziałów~\ref{chap:cpp14-type-erasure} i~\ref{chap:variadic-templates-cpp}.
\texttt{std::visit} wybiera wywołanie na podstawie aktywnego typu
\texttt{std::variant}, natomiast \texttt{std::apply} rozwija argumenty
z krotki lub podobnej struktury.

\subsection{Zestaw przeciążeń oparty na dziedziczeniu}

\begin{lstlisting}
#include <string>
#include <variant>

template <class... Fs>
struct overloaded : Fs... {
    using Fs::operator()...;
};

std::variant<int, std::string> value = std::string{"Hello"};

auto stringify = overloaded{
    [](int x) { return std::to_string(x); },
    [](const std::string& x) { return x; }
};

std::string result = std::visit(stringify, value); // "Hello"
\end{lstlisting}

Każda lambda dostarcza jeden \texttt{operator()}. Dziedziczenie
wariadyczne i deklaracja \texttt{using Fs::operator()...} tworzą zbiór
przeciążeń, spośród których kompilator wybiera odpowiednią funkcję.
W tym przykładzie wszystkie warianty wywołania zwracają ten sam typ,
\texttt{std::string}, czego wymaga tu pojedyncze wywołanie
\texttt{std::visit}. Konstrukcja jest przydatna przy przetwarzaniu
stanów programu i heterogenicznych danych.

\subsection{Rozwijanie krotki przez \texttt{std::apply}}

\begin{lstlisting}
#include <tuple>

std::tuple<int, int, int> data{2, 3, 5};
int sum = std::apply(
    [](auto... values) { return (0 + ... + values); },
    data);
// sum == 10
\end{lstlisting}

Tutaj \texttt{std::apply} bierze elementy krotki i wywołuje lambdę
jak zwykłą funkcję trójargumentową. Sama lambda wykorzystuje
\emph{fold expression}, a więc nie trzeba pisać rekurencyjnego kodu.

\section{Pakiety przechwyceń i obliczenia odroczone (C++20)}

C++20 pozwala stosować \emph{pack expansion} również w przechwytywaniu
inicjalizującym. Możemy w ten sposób zbudować jeden obiekt domknięcia,
który przechowuje cały pakiet wartości do późniejszego przetworzenia:

\begin{lstlisting}
#include <concepts>

template <std::integral... Ts>
auto delayed_sum(Ts... xs) {
    return [...stored = xs] {
        return (0 + ... + stored);
    };
}

auto job = delayed_sum(2, 3, 5);
int result = job(); // 10
\end{lstlisting}

Lambdę można zwrócić z funkcji, ponieważ wartości są przechowywane
wewnątrz domknięcia. Przechwycenie przez referencję miałoby inne
wymagania dotyczące czasu życia. W rzeczywistych fabrykach callbacków
pakiet może zawierać obiekty niekopiowalne; wtedy należy zadbać
o przekazanie własności (np. przez \texttt{std::move} i
\texttt{std::forward}) oraz o wymagania interfejsu, który callback
przechowuje.

\section{Lambdy w obliczeniach stałych}

Jeśli ciało lambdy dopuszcza obliczanie stałe, jej wywołanie może
być częścią wyrażenia stałego. Od C++17 odpowiedni operator wywołania
jest domyślnie \texttt{constexpr}, gdy spełnia wymagania takiej funkcji:

\begin{lstlisting}
constexpr auto cube = [](int x) { return x * x * x; };
static_assert(cube(3) == 27);

// C++20: each call must be evaluated at compile time.
auto fourth_power = [](int x) consteval {
    return x * x * x * x;
};
static_assert(fourth_power(2) == 16);
\end{lstlisting}

\texttt{consteval} jest mocniejszą deklaracją niż \texttt{constexpr}:
wymaga, aby właściwe wywołania funkcji natychmiastowej spełniały
reguły obliczania stałego. Nie należy traktować tego jako gwarancji,
że każda zwykła lambda zostanie wykonana podczas kompilacji.

\section{Koncepcje dla obiektów wywoływalnych}
\label{sec:callable-concepts}

W C++20 można ograniczać typy wywoływalne m.in. następującymi
koncepcjami:
\begin{itemize}
\item \texttt{std::invocable} -- poprawność wywołania z danymi argumentami;
\item \texttt{std::regular\_invocable} -- dodatkowe wymagania semantyczne;
\item \texttt{std::predicate} -- poprawnie wywoływalny predykat z
wynikiem dającym się traktować jako wartość logiczną.
\end{itemize}
Koncepcje nie dowodzą wszystkich własności semantycznych. Przykładowo
kompilator nie zweryfikuje matematycznie stabilności wyników predykatu.

\begin{lstlisting}
#include <concepts>
#include <functional>
#include <utility>

template <class F, class T>
    requires std::predicate<F&, const T&>
bool test(const T& value, F&& predicate) {
    return std::invoke(std::forward<F>(predicate), value);
}

bool ok = test(8, [](int x) { return x % 2 == 0; });
\end{lstlisting}

\textbf{Istotne rozróżnienie:} koncepcja weryfikuje przede wszystkim
własności typów i poprawność wyrażeń. Nie dowodzi matematycznie
własności predykatu, np. że comparator tworzy ścisły słaby porządek.
Programista nadal odpowiada za kontrakt semantyczny.

\section{Typowe błędy i trudniejsze przypadki}
\label{sec:lambda-pitfalls}

\subsection{Referencja przechwycona do operacji asynchronicznej}

\begin{lstlisting}
#include <future>

std::future<int> start() {
    int x = 7;
    // WRONG: the task could outlive x.
    // return std::async(std::launch::async, [&x] { return x; });
    return std::async(std::launch::async, [x] { return x; });
}
\end{lstlisting}

W zadaniach asynchronicznych preferuj przechwycenia przez wartość
lub świadomie zarządzaną własność, dopóki nie mamy ścisłej gwarancji,
że obiekt referowany przetrwa wykonanie zadania. Sama biblioteka
\texttt{std::async} nie wydłuża życia lokalnych zmiennych funkcji
uruchamiającej.

\subsection{Predykat sortowania nie może naruszać porządku}

\begin{lstlisting}
// WRONG: <= is not a strict ordering relation.
// std::sort(xs.begin(), xs.end(),
//           [](int a, int b) { return a <= b; });

std::sort(xs.begin(), xs.end(),
          [](int a, int b) { return a < b; });
\end{lstlisting}

Komparator do \texttt{sort} musi definiować \emph{ścisły słaby
porządek}. Przykładowo \texttt{comp(x,x)} ma dawać \texttt{false}.
Naruszenie wymagań algorytmu może prowadzić do niezdefiniowanego
zachowania, mimo że typy argumentów się zgadzają.

\subsection{Pomylenie kopii z oryginalnym stanem}

\begin{lstlisting}
int sum = 0;
std::for_each(xs.begin(), xs.end(), [sum](int x) mutable {
    sum += x; // modifies only the closure's copy
});
// outer sum remains 0

std::for_each(xs.begin(), xs.end(), [&sum](int x) {
    sum += x;
});
\end{lstlisting}

Nie należy zakładać, że \texttt{mutable} zmienia zewnętrzną
zmienną. Należy też pamiętać, że algorytmy
mogą kopiować obiekt funkcyjny; nie opierajmy semantyki algorytmu
na nieudokumentowanych szczegółach kopiowania predykatu.

\subsection{Przechwycenie wskaźnika nie przedłuża życia obiektu}

\begin{lstlisting}
int* p = /* pointer to some object */ nullptr;
auto f = [p] {
    if (p) return *p;
    return 0;
};
\end{lstlisting}

Wartość wskaźnika może zostać skopiowana do lambdy, ale wskazywany
obiekt nadal musi istnieć w chwili dereferencji. Przechwycenie
\texttt{shared\_ptr} przenosi własność współdzieloną, a
\texttt{weak\_ptr} pozwala ją warunkowo odzyskać. To nie jest to
samo co skopiowanie wskaźnika surowego.

\section{Zestawienie: kiedy stosować którą konstrukcję?}

\begin{center}
\begin{tabular}{|p{5.7cm}|p{6.6cm}|}
\hline
\textbf{Potrzeba} & \textbf{Zalecany wybór} \\
\hline
Krótki predykat do algorytmu & Lambda z jawnymi przechwyceniami \\
\hline
Wywołanie bez stanu w API C & Zwykła funkcja albo zgodna lambda bez przechwyceń jako wskaźnik do funkcji \\
\hline
Wydajny callback o typie znanym w czasie kompilacji & Parametr szablonowy \texttt{F}, ewentualnie \texttt{std::invocable} \\
\hline
Callback z możliwością późniejszej podmiany & \texttt{std::function} (C++20, kopiowalny cel) \\
\hline
Przechowanie \texttt{unique\_ptr} wewnątrz lambdy & Init-capture i przeniesienie; nie przechowuj takiej lambdy w \texttt{std::function} w C++20 \\
\hline
Callback żyjący dłużej niż zmienne lokalne & Kopiuj potrzebny stan lub zarządzaj życiem przez RAII i smart pointery \\
\hline
Uniwersalne wywołanie dowolnego callable & \texttt{std::invoke} i \texttt{std::forward} \\
\hline
Częściowe podstawienie kilku argumentów & Lambda lub \texttt{std::bind\_front} \\
\hline
Bezpieczny predykat dla szablonu & \texttt{std::predicate} plus kontrola jego semantyki \\
\hline
\end{tabular}
\end{center}

\paragraph{Pytania kontrolne.}
\begin{enumerate}
\item Dlaczego dwie lambdy o identycznym kodzie mają różne typy?
\item Jaka jest różnica między \texttt{[x]}, \texttt{[\&x]} i \texttt{[x]() mutable}?
\item Kiedy bezpiecznie zwrócić lambdę przechwytującą zmienną lokalną?
\item Dlaczego lambda posiadająca \texttt{unique\_ptr} nie może być celem \texttt{std::function} w C++20?
\item Jak \texttt{std::invoke} wywołuje wskaźnik do metody klasy?
\item Czym różnią się \texttt{[this]} i \texttt{[*this]}?
\item Co może pójść źle, gdy \texttt{std::views::filter} korzysta z lokalnej zmiennej przez referencję?
\item Kiedy interfejs przyjmujący szablon \texttt{F} jest lepszy od \texttt{std::function}?
\item Dlaczego komparator \texttt{a <= b} nie jest poprawny dla \texttt{std::sort}?
\item W jakiej sytuacji przechwycenie \texttt{weak\_ptr} jest lepsze od \texttt{shared\_ptr}?
\end{enumerate}

\paragraph{Dalsza lektura.}
W dokumentacji \emph{cppreference} warto przeczytać hasła
\emph{Lambda expressions}, \emph{Function objects} i \emph{Ranges}.
Dokładne opisy narzędzi znajdują się również pod hasłami:
\begin{itemize}
\item \texttt{std::invoke} oraz \texttt{std::function};
\item \texttt{std::bind\_front} oraz \texttt{std::predicate}.
\end{itemize}
